\documentclass[11pt,a4paper]{article}
\usepackage[a4paper,top=1.3cm,bottom=1.3cm,left=1.5cm,right=1.5cm,headheight=14pt,headsep=8pt,footskip=16pt]{geometry}
\usepackage{fontspec}
\setmainfont{Noto Serif CJK KR}
\setsansfont{Noto Sans CJK KR}
\setmonofont{Noto Sans Mono CJK KR}

\usepackage{amsmath,amssymb,amsfonts,bm}
\usepackage{xcolor}
\usepackage{multicol}
\usepackage{tcolorbox}
\tcbuselibrary{skins,breakable}
\usepackage{fancyhdr}
\usepackage{hyperref}

\definecolor{navy}{HTML}{002060}
\definecolor{primaryblue}{HTML}{184C78}
\definecolor{accent}{HTML}{2980B9}
\definecolor{darkgray}{HTML}{2C3E50}
\definecolor{lightbg}{HTML}{F8FAFC}
\definecolor{bordergray}{HTML}{CBD5E1}
\definecolor{successgreen}{HTML}{27AE60}

\pagestyle{fancy}
\fancyhf{}
\fancyhead[L]{\small\bfseries\color{primaryblue}가속기 인공지능 튜토리얼 \enspace$|$\enspace 제1강 슬라이드 내러티브 대본}
\fancyhead[R]{\small\color{darkgray}박종식 교수 (고려대학교)}
\fancyfoot[C]{\small\bfseries\color{primaryblue}슬라이드 \thepage\ / 35}
\fancyfoot[R]{\footnotesize\color{darkgray}AI4ACC Tutorial 2026}
\renewcommand{\headrulewidth}{0.4pt}
\renewcommand{\footrulewidth}{0.4pt}

\hypersetup{
    colorlinks=true,
    linkcolor=primaryblue,
    urlcolor=accent,
    pdftitle={Lecture 1 Slide Narratives (Korean)},
    pdfauthor={Chong Shik Park, Ph.D.}
}

\begin{document}

% --- Title Page ---
\begin{titlepage}
    \centering
    \vspace*{1.2cm}
    {\Huge\bfseries\color{navy}가속기 인공지능 응용\\[0.3em]튜토리얼 및 미니 해커톤\par}
    \vspace{0.8cm}
    {\LARGE\bfseries\color{primaryblue}제1강: AI/ML 기초 및 폐루프 궤도 제어\par}
    \vspace{0.3cm}
    {\Large\color{accent}슬라이드별 상세 구두 발표 대본 및 기술 해설서 (한국어판)\par}
    \vspace{1.2cm}

    \begin{tcolorbox}[colback=white,colframe=accent,width=0.88\textwidth,arc=4pt]
        \centering
        \textbf{강연자 및 저자:}\\[0.2em]
        {\large\bfseries 박종식 교수 (Chong Shik Park, Ph.D.)}\\[0.2em]
        고려대학교 세종캠퍼스 가속기과학과\\[0.2em]
        \texttt{kuphy@korea.ac.kr} \enspace$|$\enspace 2026년 11월 19일
    \end{tcolorbox}

    \vspace{1.2cm}
    \begin{minipage}{0.88\textwidth}
        \small\color{darkgray}
        \textbf{문서 개요 및 교수학적 구성 체계:}\\
        본 문서는 제1강 총 35개 슬라이드 전체에 대한 구두 발표 대본(Verbatim Script), 물리 이론 배경, 슬라이드 시각 구성 요약, 교육적 강조 포인트, 예상 질의응답(Q\&A), 그리고 다음 슬라이드로의 자연스러운 전환 멘트를 수록한 공식 강연 해설서입니다.
        \par\vspace{0.6em}
        각 슬라이드의 내용 전체가 정확히 1페이지 안에 완결되도록 조판(1번부터 35번 슬라이드와 1:1 페이지 대응)되어, 발표자가 강단에서 편리하게 직관적으로 슬라이드를 넘기며 강연할 수 있으며, 참가자들의 실전 해커톤 자습 교재로 최적화되었습니다.
    \end{minipage}
    \vfill
    {\footnotesize\color{darkgray}고려대학교 가속기물리연구실 \enspace$\bullet$\enspace 2026}
\end{titlepage}

% --- Table of Contents Page ---
\clearpage
\pagenumbering{roman}
\fancyfoot[C]{\small\color{darkgray}목차}

\noindent\begin{tcolorbox}[colback=navy,colframe=navy,left=8pt,right=8pt,top=4pt,bottom=4pt,boxrule=0pt,arc=2pt]
{\large\bfseries\color{white}목차 (Table of Contents) \enspace$|$\enspace 제1강 슬라이드 디렉토리 (슬라이드 1--35)}
\end{tcolorbox}
\vspace{0.4em}

{\footnotesize
\begin{multicols}{2}
\noindent\hyperlink{slide.1}{\textbf{\color{primaryblue}슬라이드 1:}\enspace 가속기 인공지능/머신러닝 기초 및 폐루프 궤도 제어\enspace\dotfill\enspace\textbf{1}} \\[0.22em]
\noindent\hyperlink{slide.2}{\textbf{\color{primaryblue}슬라이드 2:}\enspace 전체 교육과정 및 마스터 프로그램 일정\enspace\dotfill\enspace\textbf{2}} \\[0.22em]
\noindent\hyperlink{slide.3}{\textbf{\color{primaryblue}슬라이드 3:}\enspace 현대 입자 가속기에서 왜 인공지능인가?\enspace\dotfill\enspace\textbf{3}} \\[0.22em]
\noindent\hyperlink{slide.4}{\textbf{\color{primaryblue}슬라이드 4:}\enspace 가속기 인공지능 핵심 분야 (1): 대리 모델 및 자율 최적화\enspace\dotfill\enspace\textbf{4}} \\[0.22em]
\noindent\hyperlink{slide.5}{\textbf{\color{primaryblue}슬라이드 5:}\enspace 가속기 인공지능 핵심 분야 (2): 진단, 장비 보호 및 자율 에이전트\enspace\dotfill\enspace\textbf{5}} \\[0.22em]
\noindent\hyperlink{slide.6}{\textbf{\color{primaryblue}슬라이드 6:}\enspace 힐 방정식(Hill's Equation)과 횡방향 입자 운동\enspace\dotfill\enspace\textbf{6}} \\[0.22em]
\noindent\hyperlink{slide.7}{\textbf{\color{primaryblue}슬라이드 7:}\enspace 베타트론 운동과 쿠랑-스나이더(Courant-Snyder) 광학\enspace\dotfill\enspace\textbf{7}} \\[0.22em]
\noindent\hyperlink{slide.8}{\textbf{\color{primaryblue}슬라이드 8:}\enspace 15미터 빔 전송 격자 구조 및 진단 장치 배치\enspace\dotfill\enspace\textbf{8}} \\[0.22em]
\noindent\hyperlink{slide.9}{\textbf{\color{primaryblue}슬라이드 9:}\enspace 궤도 응답 행렬(Orbit Response Matrix, ORM) 유도\enspace\dotfill\enspace\textbf{9}} \\[0.22em]
\noindent\hyperlink{slide.10}{\textbf{\color{primaryblue}슬라이드 10:}\enspace 하드웨어 구동 한계 및 센서 노이즈 제약\enspace\dotfill\enspace\textbf{10}} \\[0.22em]
\noindent\hyperlink{slide.11}{\textbf{\color{primaryblue}슬라이드 11:}\enspace 전자석 철심의 비선형 자기 포화(Magnetic Core Saturation)\enspace\dotfill\enspace\textbf{11}} \\[0.22em]
\noindent\hyperlink{slide.12}{\textbf{\color{primaryblue}슬라이드 12:}\enspace 인공 뉴런(Artificial Neuron)과 신경망 계층 구성\enspace\dotfill\enspace\textbf{12}} \\[0.22em]
\noindent\hyperlink{slide.13}{\textbf{\color{primaryblue}슬라이드 13:}\enspace 신경망 활성화 함수(Activation Functions)의 역할과 분류\enspace\dotfill\enspace\textbf{13}} \\[0.22em]
\noindent\hyperlink{slide.14}{\textbf{\color{primaryblue}슬라이드 14:}\enspace 신경망 대리 모델 예제: 입자 전달 행렬(Transfer Matrix M) 학습\enspace\dotfill\enspace\textbf{14}} \\[0.22em]
\noindent\hyperlink{slide.15}{\textbf{\color{primaryblue}슬라이드 15:}\enspace 빔 전송 대리 모델에서 활성화 함수별 위상 공간 변형\enspace\dotfill\enspace\textbf{15}} \\[0.22em]
\noindent\hyperlink{slide.16}{\textbf{\color{primaryblue}슬라이드 16:}\enspace 머신러닝 최적화 알고리즘: 경사하강법부터 Adam 및 L-BFGS-B까지\enspace\dotfill\enspace\textbf{16}} \\[0.22em]
\noindent\hyperlink{slide.17}{\textbf{\color{primaryblue}슬라이드 17:}\enspace MLP를 넘어선 가속기 특화 신경망 구조: PINN, PRNN, VAE, U-Net\enspace\dotfill\enspace\textbf{17}} \\[0.22em]
\noindent\hyperlink{slide.18}{\textbf{\color{primaryblue}슬라이드 18:}\enspace 지도학습: 디지털 트윈 대리 모델링 및 하이퍼파라미터\enspace\dotfill\enspace\textbf{18}} \\[0.22em]
\noindent\hyperlink{slide.19}{\textbf{\color{primaryblue}슬라이드 19:}\enspace 신경망 대리 모델 디지털 트윈 구조\enspace\dotfill\enspace\textbf{19}} \\[0.22em]
\noindent\hyperlink{slide.20}{\textbf{\color{primaryblue}슬라이드 20:}\enspace 대리 모델 훈련 및 독립 홀드아웃(Holdout) 검증 결과\enspace\dotfill\enspace\textbf{20}} \\[0.22em]
\noindent\hyperlink{slide.21}{\textbf{\color{primaryblue}슬라이드 21:}\enspace 폐루프 제어 아키텍처 및 대리 모델 역산 흐름\enspace\dotfill\enspace\textbf{21}} \\[0.22em]
\noindent\hyperlink{slide.22}{\textbf{\color{primaryblue}슬라이드 22:}\enspace 베이스라인 0: 고전적 선형 SVD 의사역행렬 제어기\enspace\dotfill\enspace\textbf{22}} \\[0.22em]
\noindent\hyperlink{slide.23}{\textbf{\color{primaryblue}슬라이드 23:}\enspace 베이스라인 1 (트랙 1): 딥러닝 대리 모델 역산 최적화\enspace\dotfill\enspace\textbf{23}} \\[0.22em]
\noindent\hyperlink{slide.24}{\textbf{\color{primaryblue}슬라이드 24:}\enspace 강화학습: 마르코프 결정 과정(MDP) 정식화 및 하이퍼파라미터\enspace\dotfill\enspace\textbf{24}} \\[0.22em]
\noindent\hyperlink{slide.25}{\textbf{\color{primaryblue}슬라이드 25:}\enspace 트랙 2: 물리 사전 지식 결합 폐루프 강화학습 정책\enspace\dotfill\enspace\textbf{25}} \\[0.22em]
\noindent\hyperlink{slide.26}{\textbf{\color{primaryblue}슬라이드 26:}\enspace 동적 조화 스텝 감쇠(Dynamic Harmonic Step Damping)\enspace\dotfill\enspace\textbf{26}} \\[0.22em]
\noindent\hyperlink{slide.27}{\textbf{\color{primaryblue}슬라이드 27:}\enspace 베이지안 최적화: 확률론적 가우시안 프로세스(GP) 및 하이퍼파라미터\enspace\dotfill\enspace\textbf{27}} \\[0.22em]
\noindent\hyperlink{slide.28}{\textbf{\color{primaryblue}슬라이드 28:}\enspace 트랙 3: 가우시안 프로세스 기반 폐루프 베이지안 최적화\enspace\dotfill\enspace\textbf{28}} \\[0.22em]
\noindent\hyperlink{slide.29}{\textbf{\color{primaryblue}슬라이드 29:}\enspace 하한 신뢰 구간(Lower Confidence Bound, LCB) 획득 함수\enspace\dotfill\enspace\textbf{29}} \\[0.22em]
\noindent\hyperlink{slide.30}{\textbf{\color{primaryblue}슬라이드 30:}\enspace 체육관(Gymnasium) 환경 API 표준 규격\enspace\dotfill\enspace\textbf{30}} \\[0.22em]
\noindent\hyperlink{slide.31}{\textbf{\color{primaryblue}슬라이드 31:}\enspace 공식 30개 시드 평가 프로토콜 및 채점 공식\enspace\dotfill\enspace\textbf{31}} \\[0.22em]
\noindent\hyperlink{slide.32}{\textbf{\color{primaryblue}슬라이드 32:}\enspace 공식 벤치마크 리더보드 및 수렴 궤적 비교\enspace\dotfill\enspace\textbf{32}} \\[0.22em]
\noindent\hyperlink{slide.33}{\textbf{\color{primaryblue}슬라이드 33:}\enspace 제어 방식별 종합 비교 요약 (기법, 연산 지연 및 성능)\enspace\dotfill\enspace\textbf{33}} \\[0.22em]
\noindent\hyperlink{slide.34}{\textbf{\color{primaryblue}슬라이드 34:}\enspace 해커톤 참가자를 위한 실전 공략 전략 (TODO 1, 2, 3)\enspace\dotfill\enspace\textbf{34}} \\[0.22em]
\noindent\hyperlink{slide.35}{\textbf{\color{primaryblue}슬라이드 35:}\enspace 강의 요약, 핵심 교훈 및 제2강 연계 안내\enspace\dotfill\enspace\textbf{35}} \\[0.22em]
\end{multicols}}

% --- Slide Pages (1 to 35) ---
\clearpage
\pagenumbering{arabic}
\setcounter{page}{1}
\fancyfoot[C]{\small\bfseries\color{primaryblue}슬라이드 \thepage\ / 35}

\clearpage
\hypertarget{slide.1}{}
\phantomsection
\addcontentsline{toc}{section}{슬라이드 1: 가속기 인공지능/머신러닝 기초 및 폐루프 궤도 제어}
\noindent\begin{tcolorbox}[colback=navy,colframe=navy,left=6pt,right=6pt,top=3pt,bottom=3pt,boxrule=0pt,arc=2pt]
{\large\bfseries\color{white}슬라이드 1:\enspace 가속기 인공지능/머신러닝 기초 및 폐루프 궤도 제어}\hfill{\footnotesize\bfseries\color{white!85!lightbg}Title Slide \enspace$|$\enspace 슬라이드 1 / 35}
\end{tcolorbox}
\vspace{0.2em}
\begin{tcolorbox}[colback=lightbg,colframe=bordergray,title={\footnotesize\bfseries\color{darkgray}슬라이드 시각 구성 요약},fontupper=\small,left=6pt,right=6pt,top=2.5pt,bottom=2.5pt,boxrule=0.6pt,arc=2.5pt]
고려대학교 로고, 강좌 제목, 연사 정보(박종식 교수), 해커톤 주제가 포함된 타이틀 슬라이드.
\end{tcolorbox}
\vspace{0.2em}
\begin{tcolorbox}[colback=white,colframe=primaryblue,title={\footnotesize\bfseries\color{white}발표자 구두 대본 (Spoken Script)},fontupper=\small,left=6pt,right=6pt,top=3pt,bottom=3pt,boxrule=0.8pt,arc=2.5pt]
안녕하십니까, 참가자 여러분. 가속기 인공지능 응용 튜토리얼 및 미니 해커톤의 첫 번째 강의에 오신 것을 환영합니다. 저는 고려대학교 세종캠퍼스 가속기과학과의 박종식입니다. 오늘 첫 번째 강의에서는 고전 가속기 빔 물리와 현대 인공지능 기법을 체계적으로 연결하고자 합니다. 횡방향 빔 동역학, 힐 방정식(Hill's equation), 궤도 응답 행렬(ORM)과 같은 기본 물리 원리에서 출발하여, 신경망 기반 디지털 트윈 대리 모델(surrogate model), 물리 사전 지식 기반 강화학습(RL), 그리고 베이지안 최적화(BO)에 이르는 폐루프 제어 기법을 심도 있게 다룰 것입니다. 오늘 강의의 목표는 단순한 이론 학습을 넘어, 오후에 진행될 미니 해커톤에서 여러분 스스로 고성능 자율 제어기를 직접 설계하고 벤치마크할 수 있는 실질적인 역량을 갖추는 것입니다.
\end{tcolorbox}
\vspace{0.2em}
\begin{tcolorbox}[colback=lightbg,colframe=bordergray,left=6pt,right=6pt,top=3pt,bottom=3pt,boxrule=0.6pt,arc=2.5pt]
\small
\noindent\textbf{\color{navy}교수학적 강조점 (Focus Point):}\enspace 강의의 실무적·학술적 가치를 강조하고, 물리학적 기본 원리와 최신 인공지능 알고리즘이 상호 보완적으로 작동함을 밝힙니다.\par\vspace{0.3em}
\noindent\textbf{\color{successgreen}예상 질문 및 모범 답변 (Q\&A):}\enspace Q: 딥러닝 사전 경험이 부족해도 해커톤 참여가 가능한가요? A: 네, 선형 전달 행렬 기초부터 단계별로 설명하므로 충분히 따라오실 수 있습니다.\par\vspace{0.3em}
\noindent\textbf{\color{accent}다음 슬라이드 연결 멘트 (Transition):}\enspace 먼저 오늘 하루의 전체 일정을 살펴보겠습니다.
\end{tcolorbox}
\clearpage
\hypertarget{slide.2}{}
\phantomsection
\addcontentsline{toc}{section}{슬라이드 2: 전체 교육과정 및 마스터 프로그램 일정}
\noindent\begin{tcolorbox}[colback=navy,colframe=navy,left=6pt,right=6pt,top=3pt,bottom=3pt,boxrule=0pt,arc=2pt]
{\large\bfseries\color{white}슬라이드 2:\enspace 전체 교육과정 및 마스터 프로그램 일정}\hfill{\footnotesize\bfseries\color{white!85!lightbg}Workshop Curriculum \enspace$|$\enspace 슬라이드 2 / 35}
\end{tcolorbox}
\vspace{0.2em}
\begin{tcolorbox}[colback=lightbg,colframe=bordergray,title={\footnotesize\bfseries\color{darkgray}슬라이드 시각 구성 요약},fontupper=\small,left=6pt,right=6pt,top=2.5pt,bottom=2.5pt,boxrule=0.6pt,arc=2.5pt]
3단계 일정 카드: 제1강(기초 이론 및 제어), 제2강(Thorsten Hellert 박사의 에이전틱 AI), 오후 미니 해커톤 실습 및 실시간 리더보드.
\end{tcolorbox}
\vspace{0.2em}
\begin{tcolorbox}[colback=white,colframe=primaryblue,title={\footnotesize\bfseries\color{white}발표자 구두 대본 (Spoken Script)},fontupper=\small,left=6pt,right=6pt,top=3pt,bottom=3pt,boxrule=0.8pt,arc=2.5pt]
오늘 프로그램은 체계적인 이론 교육에서 시작하여 참가자 주도의 실전 해커톤으로 이어집니다. 오전 제1강에서는 제가 횡방향 빔 물리와 대리 모델링, 세 가지 머신러닝 패러다임을 강의합니다. 이어지는 제2강에서는 미국 로렌스 버클리 국립연구소(LBNL)의 Thorsten Hellert 박사님께서 최신 에이전틱 AI(Agentic AI) 기술을 소개합니다. 거대언어모델(LLM)이 파이썬 API 도구를 직접 실행하며 가속기 시뮬레이션을 자율 오케스트레이션하는 기법을 배울 수 있습니다. 그리고 오후 1시 30분부터는 본격적인 미니 해커톤이 3시간 30분 동안 펼쳐집니다. 15미터 빔 전송선 시뮬레이션 환경에서 여러분만의 제어기를 튜닝하고, Weights \& Biases 실시간 리더보드를 통해 30개 표준 시드에 대한 최종 점수를 경쟁하게 됩니다.
\end{tcolorbox}
\vspace{0.2em}
\begin{tcolorbox}[colback=lightbg,colframe=bordergray,left=6pt,right=6pt,top=3pt,bottom=3pt,boxrule=0.6pt,arc=2.5pt]
\small
\noindent\textbf{\color{navy}교수학적 강조점 (Focus Point):}\enspace 오전 강의에서 배우는 이론들이 오후 해커톤 과제(TODO 1, 2, 3)에 직접 어떻게 활용되는지 연계성을 강조합니다.\par\vspace{0.3em}
\noindent\textbf{\color{successgreen}예상 질문 및 모범 답변 (Q\&A):}\enspace Q: 해커톤 평가는 어떻게 이루어지나요? A: 30개 표준 시드에 대해 성공 여부, 수렴 속도, 안착 정밀도를 종합한 1000점 만점 공식으로 평가합니다.\par\vspace{0.3em}
\noindent\textbf{\color{accent}다음 슬라이드 연결 멘트 (Transition):}\enspace 본격적인 내용에 앞서, 현대 가속기 시설에서 왜 인공지능 기법이 필요한지 그 배경을 살펴보겠습니다.
\end{tcolorbox}
\clearpage
\hypertarget{slide.3}{}
\phantomsection
\addcontentsline{toc}{section}{슬라이드 3: 현대 입자 가속기에서 왜 인공지능인가?}
\noindent\begin{tcolorbox}[colback=navy,colframe=navy,left=6pt,right=6pt,top=3pt,bottom=3pt,boxrule=0pt,arc=2pt]
{\large\bfseries\color{white}슬라이드 3:\enspace 현대 입자 가속기에서 왜 인공지능인가?}\hfill{\footnotesize\bfseries\color{white!85!lightbg}Operational Context \enspace$|$\enspace 슬라이드 3 / 35}
\end{tcolorbox}
\vspace{0.2em}
\begin{tcolorbox}[colback=lightbg,colframe=bordergray,title={\footnotesize\bfseries\color{darkgray}슬라이드 시각 구성 요약},fontupper=\small,left=6pt,right=6pt,top=2.5pt,bottom=2.5pt,boxrule=0.6pt,arc=2.5pt]
고강도 빔 안정성 요건, 서브미크론 허용오차, 빔 손실로 인한 방사화 및 초전도 챔버 파손 위험, 운전원의 인지 한계를 정리한 개요 카드.
\end{tcolorbox}
\vspace{0.2em}
\begin{tcolorbox}[colback=white,colframe=primaryblue,title={\footnotesize\bfseries\color{white}발표자 구두 대본 (Spoken Script)},fontupper=\small,left=6pt,right=6pt,top=3pt,bottom=3pt,boxrule=0.8pt,arc=2.5pt]
4세대 방사광가속기, 자유전자레이저(FEL), 고강도 중성자 파쇄원 등 현대 가속기는 인류가 제작한 가장 정밀하고 복잡한 장치 중 하나입니다. 이러한 최첨단 시설에서는 횡방향 빔 궤도를 수백 나노미터에서 서브미크론 수준으로 일정하게 유지해야 합니다. 동시에 메가와트급 빔 에너지를 다루기 때문에, 단 수 마이크로초 동안의 미세한 빔 손실(beam loss)만으로도 진공 챔버가 손상되거나 초전도 가속관(SRF)이 퀜치(quench)되어 수억 원 이상의 운영 손실과 수주일간의 빔 정지를 초래할 수 있습니다. 기존의 비례-적분(PI) 및 선형 피드백 제어기는 가속기 시스템이 선형이고 시불변이라는 가정에 의존합니다. 하지만 실제 가속기 전자석은 자기 포화, 히스테리시스, 열적 드리프트 등 강한 비선형성을 지닙니다. 수백 개의 제어 변수를 인간 운전원이 수동으로 최적화하는 것은 한계에 도달했습니다. 따라서 고속 대리 모델, 자율 다목적 최적화, 밀리초 단위 이상 진단을 수행하는 인공지능 도입이 필연적입니다.
\end{tcolorbox}
\vspace{0.2em}
\begin{tcolorbox}[colback=lightbg,colframe=bordergray,left=6pt,right=6pt,top=3pt,bottom=3pt,boxrule=0.6pt,arc=2.5pt]
\small
\noindent\textbf{\color{navy}교수학적 강조점 (Focus Point):}\enspace 가속기 운전 현장의 실제 위험성(장비 파손, 빔 품질 저하, 빔 다운타임)을 강조하여 AI 도입의 필요성을 각인시킵니다.\par\vspace{0.3em}
\noindent\textbf{\color{successgreen}예상 질문 및 모범 답변 (Q\&A):}\enspace Q: 기존의 PID 제어로 비선형 드리프트를 해결할 수 없나요? A: 기준점 부근의 미세 선형 구간에서는 가능하지만, 포화나 외란 발생 시 발진하거나 수렴 속도가 급격히 저하됩니다.\par\vspace{0.3em}
\noindent\textbf{\color{accent}다음 슬라이드 연결 멘트 (Transition):}\enspace 그렇다면 가속기 과학 전반에서 인공지능은 어떠한 핵심 기둥들로 분류되어 적용되고 있는지 살펴보겠습니다.
\end{tcolorbox}
\clearpage
\hypertarget{slide.4}{}
\phantomsection
\addcontentsline{toc}{section}{슬라이드 4: 가속기 인공지능 핵심 분야 (1): 대리 모델 및 자율 최적화}
\noindent\begin{tcolorbox}[colback=navy,colframe=navy,left=6pt,right=6pt,top=3pt,bottom=3pt,boxrule=0pt,arc=2pt]
{\large\bfseries\color{white}슬라이드 4:\enspace 가속기 인공지능 핵심 분야 (1): 대리 모델 및 자율 최적화}\hfill{\footnotesize\bfseries\color{white!85!lightbg}Domain Overview \enspace$|$\enspace 슬라이드 4 / 35}
\end{tcolorbox}
\vspace{0.2em}
\begin{tcolorbox}[colback=lightbg,colframe=bordergray,title={\footnotesize\bfseries\color{darkgray}슬라이드 시각 구성 요약},fontupper=\small,left=6pt,right=6pt,top=2.5pt,bottom=2.5pt,boxrule=0.6pt,arc=2.5pt]
핵심 기둥 1: 고속 디지털 트윈 대리 모델(초고속 시뮬레이션), 핵심 기둥 2: 가우시안 프로세스 베이지안 최적화를 통한 자율 다목적 튜닝.
\end{tcolorbox}
\vspace{0.2em}
\begin{tcolorbox}[colback=white,colframe=primaryblue,title={\footnotesize\bfseries\color{white}발표자 구두 대본 (Spoken Script)},fontupper=\small,left=6pt,right=6pt,top=3pt,bottom=3pt,boxrule=0.8pt,arc=2.5pt]
가속기 인공지능 응용은 흔히 다섯 가지 기둥으로 나뉩니다. 이 슬라이드에서는 첫 번째와 두 번째 기둥을 살펴봅니다. 첫 번째 기둥은 '신경망 대리 모델(Digital Twin Surrogates)'입니다. ELEGANT나 MAD-X 같은 정밀 수치 시뮬레이션 코드는 공간전하 효과와 집단 현상을 정확히 모사하지만, 1회 계산에 수 분에서 수 시간이 소요됩니다. 딥러닝 대리 모델은 이러한 비선형 동역학을 사전에 학습하여 0.2 밀리초 이하의 초고속으로 결과를 예측합니다. 이를 통해 실시간 모델 예측 제어(MPC)가 가능해집니다. 두 번째 기둥은 '자율 다목적 최적화(Autonomous Multi-Objective Optimization)'입니다. 가속기 운전 시 빔 에미턴스 최소화와 빔 투과율 최대화는 종종 상충(trade-off) 관계에 놓입니다. 가우시안 프로세스(GP) 기반 베이지안 최적화는 불확실성을 모델링하여 수천 번의 무작위 스캔 대신 수십 회의 최소 빔 측정만으로 최적의 파레토 프론트(Pareto front)를 찾아냅니다.
\end{tcolorbox}
\vspace{0.2em}
\begin{tcolorbox}[colback=lightbg,colframe=bordergray,left=6pt,right=6pt,top=3pt,bottom=3pt,boxrule=0.6pt,arc=2.5pt]
\small
\noindent\textbf{\color{navy}교수학적 강조점 (Focus Point):}\enspace 정밀 시뮬레이션 코드의 느린 계산 시간과 딥러닝 대리 모델의 밀리초 이하 추론 속도를 명확히 대조합니다.\par\vspace{0.3em}
\noindent\textbf{\color{successgreen}예상 질문 및 모범 답변 (Q\&A):}\enspace Q: 대리 모델이 학습 범위를 벗어난 영역에서도 정확한가요? A: 순수 MLP는 외삽 능력이 떨어지므로, 물리 사전 지식이나 적절한 규제 기법이 필수적입니다.\par\vspace{0.3em}
\noindent\textbf{\color{accent}다음 슬라이드 연결 멘트 (Transition):}\enspace 이어서 가속기 AI의 나머지 세 기둥인 가상 진단, 장비 보호, 에이전틱 AI를 살펴보겠습니다.
\end{tcolorbox}
\clearpage
\hypertarget{slide.5}{}
\phantomsection
\addcontentsline{toc}{section}{슬라이드 5: 가속기 인공지능 핵심 분야 (2): 진단, 장비 보호 및 자율 에이전트}
\noindent\begin{tcolorbox}[colback=navy,colframe=navy,left=6pt,right=6pt,top=3pt,bottom=3pt,boxrule=0pt,arc=2pt]
{\large\bfseries\color{white}슬라이드 5:\enspace 가속기 인공지능 핵심 분야 (2): 진단, 장비 보호 및 자율 에이전트}\hfill{\footnotesize\bfseries\color{white!85!lightbg}Domain Overview \enspace$|$\enspace 슬라이드 5 / 35}
\end{tcolorbox}
\vspace{0.2em}
\begin{tcolorbox}[colback=lightbg,colframe=bordergray,title={\footnotesize\bfseries\color{darkgray}슬라이드 시각 구성 요약},fontupper=\small,left=6pt,right=6pt,top=2.5pt,bottom=2.5pt,boxrule=0.6pt,arc=2.5pt]
핵심 기둥 3: 가상 진단 및 비파괴 위상공간 복원, 핵심 기둥 4: 이상 감지, 고장 예측 및 장비 보호, 핵심 기둥 5: 에이전틱 AI 및 자율 운전 워크플로.
\end{tcolorbox}
\vspace{0.2em}
\begin{tcolorbox}[colback=white,colframe=primaryblue,title={\footnotesize\bfseries\color{white}발표자 구두 대본 (Spoken Script)},fontupper=\small,left=6pt,right=6pt,top=3pt,bottom=3pt,boxrule=0.8pt,arc=2.5pt]
세 번째 기둥은 '가상 진단(virtual diagnostics)'입니다. 종방향 번치 프로파일이나 슬라이스 에미턴스 같은 빔 특성은 스크린, 와이어 스캐너, 편향 공동처럼 빔을 교란하거나 파괴하는 장치로만 측정할 수 있습니다. 가상 진단은 RF 신호나 BPM 측정값 같은 비파괴 신호로부터 이러한 물리량을 샷 단위로 예측하는 지도학습 모델입니다. 네 번째 기둥은 '이상 감지, 고장 예측 및 장비 보호'입니다. 비지도 오토인코더와 순환신경망이 수천 개의 텔레메트리 채널을 감시하여 전조 드리프트나 RF 브레이크다운, 공동 고장의 징후를 찾아내고, 운전원과 인터록이 조기에 대응하여 가동 중단 시간을 줄일 수 있게 합니다. 다섯 번째 기둥은 가장 새로운 분야인 '에이전틱 AI(Agentic AI)'입니다. 거대언어모델(LLM)을 파이썬 실행 환경, 시뮬레이션 코드, 제어 시스템 인터페이스와 연결하여, 에이전트가 진단 정보를 읽고 다단계 절차를 계획한 뒤 사람의 검토와 안전 가드레일 아래에서 수행합니다.
\end{tcolorbox}
\vspace{0.2em}
\begin{tcolorbox}[colback=lightbg,colframe=bordergray,left=6pt,right=6pt,top=3pt,bottom=3pt,boxrule=0.6pt,arc=2.5pt]
\small
\noindent\textbf{\color{navy}교수학적 강조점 (Focus Point):}\enspace 다섯 번째 기둥인 에이전틱 AI는 제2강에서 Thorsten Hellert 박사가 구체적인 실무 사례와 함께 집중적으로 다룰 것임을 안내합니다.\par\vspace{0.3em}
\noindent\textbf{\color{successgreen}예상 질문 및 모범 답변 (Q\&A):}\enspace Q: 가상 진단이란 무엇인가요? A: 빔을 차단해야만 측정할 수 있는 물리량을, 비파괴 센서 측정값으로부터 머신러닝으로 실시간 추론해내는 기술입니다.\par\vspace{0.3em}
\noindent\textbf{\color{accent}다음 슬라이드 연결 멘트 (Transition):}\enspace 이제 가속기 인공지능의 물리적 토대가 되는 횡방향 빔 동역학의 지배 방정식으로 들어가 보겠습니다.
\end{tcolorbox}
\clearpage
\hypertarget{slide.6}{}
\phantomsection
\addcontentsline{toc}{section}{슬라이드 6: 힐 방정식(Hill's Equation)과 횡방향 입자 운동}
\noindent\begin{tcolorbox}[colback=navy,colframe=navy,left=6pt,right=6pt,top=3pt,bottom=3pt,boxrule=0pt,arc=2pt]
{\large\bfseries\color{white}슬라이드 6:\enspace 힐 방정식(Hill's Equation)과 횡방향 입자 운동}\hfill{\footnotesize\bfseries\color{white!85!lightbg}Physics Foundations \enspace$|$\enspace 슬라이드 6 / 35}
\end{tcolorbox}
\vspace{0.2em}
\begin{tcolorbox}[colback=lightbg,colframe=bordergray,title={\footnotesize\bfseries\color{darkgray}슬라이드 시각 구성 요약},fontupper=\small,left=6pt,right=6pt,top=2.5pt,bottom=2.5pt,boxrule=0.6pt,arc=2.5pt]
로렌츠 힘 방정식, 2극/4극 자장 선형 전개, 집속 구배 K(s), 주기적 복원력을 나타내는 힐의 2계 선형 미분방정식.
\end{tcolorbox}
\vspace{0.2em}
\begin{tcolorbox}[colback=white,colframe=primaryblue,title={\footnotesize\bfseries\color{white}발표자 구두 대본 (Spoken Script)},fontupper=\small,left=6pt,right=6pt,top=3pt,bottom=3pt,boxrule=0.8pt,arc=2.5pt]
횡방향 입자 운동은 로렌츠 힘 방정식에서 시작됩니다: d p / d t = q ( E + v × B ). 상대론적 속도를 갖는 고에너지 빔에서는 자기장에 의한 편향력이 전기장보다 광속 c배만큼 강력하므로, 빔의 궤도 수정과 집속에는 전자석이 주로 사용됩니다. 설계 궤도 부근에서 수직 자장 B\_y를 테일러 전개하면 2극 자장 B\_0, 4극 자장 구배 G = ∂B\_y/∂x, 그리고 6극 이상의 고차 다극 자장이 나타납니다. 자장 구배 G를 빔의 자기 강도(magnetic rigidity) B-rho로 정규화하면 집속 함수 K(s) = q G / p가 정의됩니다. 곡선 좌표계에서 운동방정식을 정리하면 바로 가속기 물리학의 근간인 '힐 방정식(Hill's equation)': x''(s) + K(s)x(s) = 0을 얻게 됩니다. 이는 용수철 상수 K(s)가 가속기 위치 s에 따라 주기적으로 변하는 조화 진동자(harmonic oscillator) 방정식과 수학적으로 완전히 동일합니다.
\end{tcolorbox}
\vspace{0.2em}
\begin{tcolorbox}[colback=lightbg,colframe=bordergray,left=6pt,right=6pt,top=3pt,bottom=3pt,boxrule=0.6pt,arc=2.5pt]
\small
\noindent\textbf{\color{navy}교수학적 강조점 (Focus Point):}\enspace K(s) > 0이면 수평 집속(QF), K(s) < 0이면 수평 발산(QD)이 일어나며, 빔이 중심 궤도를 따라 주기적 진동을 수행함을 설명합니다.\par\vspace{0.3em}
\noindent\textbf{\color{successgreen}예상 질문 및 모범 답변 (Q\&A):}\enspace Q: 왜 K가 위치 s의 함수인가요? A: 빔이 진행하면서 집속 전자석, 발산 전자석, 자장이 없는 드리프트 구간을 번갈아 통과하기 때문입니다.\par\vspace{0.3em}
\noindent\textbf{\color{accent}다음 슬라이드 연결 멘트 (Transition):}\enspace 위치 종속적 복원력을 지닌 힐 방정식을 해석적으로 풀기 위해 쿠랑-스나이더(Courant-Snyder) 광학 변수를 도입합니다.
\end{tcolorbox}
\clearpage
\hypertarget{slide.7}{}
\phantomsection
\addcontentsline{toc}{section}{슬라이드 7: 베타트론 운동과 쿠랑-스나이더(Courant-Snyder) 광학}
\noindent\begin{tcolorbox}[colback=navy,colframe=navy,left=6pt,right=6pt,top=3pt,bottom=3pt,boxrule=0pt,arc=2pt]
{\large\bfseries\color{white}슬라이드 7:\enspace 베타트론 운동과 쿠랑-스나이더(Courant-Snyder) 광학}\hfill{\footnotesize\bfseries\color{white!85!lightbg}Beam Optics \enspace$|$\enspace 슬라이드 7 / 35}
\end{tcolorbox}
\vspace{0.2em}
\begin{tcolorbox}[colback=lightbg,colframe=bordergray,title={\footnotesize\bfseries\color{darkgray}슬라이드 시각 구성 요약},fontupper=\small,left=6pt,right=6pt,top=2.5pt,bottom=2.5pt,boxrule=0.6pt,arc=2.5pt]
플로케 해 x(s) = √(ε·β(s))·cos(ψ(s)+$\phi_0$), 베타 함수 β(s), 위상 전진 ψ(s), 에미턴스 ε 및 빔 포락선(envelope) 다이어그램.
\end{tcolorbox}
\vspace{0.2em}
\begin{tcolorbox}[colback=white,colframe=primaryblue,title={\footnotesize\bfseries\color{white}발표자 구두 대본 (Spoken Script)},fontupper=\small,left=6pt,right=6pt,top=3pt,bottom=3pt,boxrule=0.8pt,arc=2.5pt]
K(s)가 위치에 따라 변하므로, 힐 방정식의 해는 단순한 고정 진폭 사인파가 될 수 없습니다. 1958년 쿠랑과 스나이더는 플로케(Floquet) 정리를 적용하여 일반해가 x(s) = √(ε·β(s))·cos(ψ(s) + $\phi_0$)의 형태로 표현됨을 증명했습니다. 여기서 β(s)는 '베타트론 진폭 함수(beta function)'로, 국소 빔 크기의 포락선(envelope)을 결정합니다. β(s)가 큰 곳에서는 빔 폭이 넓어지고, β(s)가 작은 곳에서는 빔이 강하게 집속됩니다. 위상 전진 ψ(s)는 1/β를 적분한 값으로, 입자가 진동 위상을 얼마나 빠르게 축적하는지 나타냅니다. 에미턴스 ε은 빔이 위상 공간에서 차지하는 불변 면적입니다. α, β, γ로 구성되는 트위스(Twiss) 매개변수는 선형 빔 광학 계산의 핵심 언어입니다.
\end{tcolorbox}
\vspace{0.2em}
\begin{tcolorbox}[colback=lightbg,colframe=bordergray,left=6pt,right=6pt,top=3pt,bottom=3pt,boxrule=0.6pt,arc=2.5pt]
\small
\noindent\textbf{\color{navy}교수학적 강조점 (Focus Point):}\enspace 실제 빔 폭 표준편차 σ(s) = √(ε·β(s)) 관계를 언급하여 물리적인 빔 크기와 베타 함수의 직접적 연결고리를 설명합니다.\par\vspace{0.3em}
\noindent\textbf{\color{successgreen}예상 질문 및 모범 답변 (Q\&A):}\enspace Q: 가속 과정에서 에미턴스는 보존되나요? A: 기하학적 에미턴스는 단열 감쇠(adiabatic damping)로 줄어들지만, 정규화 에미턴스는 보존됩니다.\par\vspace{0.3em}
\noindent\textbf{\color{accent}다음 슬라이드 연결 멘트 (Transition):}\enspace 보정 전자석이 이 진동에 어떻게 작용하는지 유도하기 전에, 오늘 오후에 제어하게 될 빔라인을 먼저 살펴보겠습니다.
\end{tcolorbox}
\clearpage
\hypertarget{slide.8}{}
\phantomsection
\addcontentsline{toc}{section}{슬라이드 8: 15미터 빔 전송 격자 구조 및 진단 장치 배치}
\noindent\begin{tcolorbox}[colback=navy,colframe=navy,left=6pt,right=6pt,top=3pt,bottom=3pt,boxrule=0pt,arc=2pt]
{\large\bfseries\color{white}슬라이드 8:\enspace 15미터 빔 전송 격자 구조 및 진단 장치 배치}\hfill{\footnotesize\bfseries\color{white!85!lightbg}Hackathon Benchmark \enspace$|$\enspace 슬라이드 8 / 35}
\end{tcolorbox}
\vspace{0.2em}
\begin{tcolorbox}[colback=lightbg,colframe=bordergray,title={\footnotesize\bfseries\color{darkgray}슬라이드 시각 구성 요약},fontupper=\small,left=6pt,right=6pt,top=2.5pt,bottom=2.5pt,boxrule=0.6pt,arc=2.5pt]
15.0m 전송선 모식도: 교대 배치된 6개의 4극자(QF/QD), 6개의 2극 보정 전자석(HCM1~3, VCM1~3), 12개의 BPM 진단 스테이션.
\end{tcolorbox}
\vspace{0.2em}
\begin{tcolorbox}[colback=white,colframe=primaryblue,title={\footnotesize\bfseries\color{white}발표자 구두 대본 (Spoken Script)},fontupper=\small,left=6pt,right=6pt,top=3pt,bottom=3pt,boxrule=0.8pt,arc=2.5pt]
이 슬라이드는 오늘 해커톤에서 다룰 15.0미터 길이의 빔 전송선 물리 배치도입니다. 교대 배치된 6개의 4극 자석(QF1~3, QD1~3)이 FODO와 유사한 교대 구배 집속 격자를 구성하고 있습니다. 궤도 제어를 위한 구동기(actuator)로는 수평 스티어러 3대(HCM1, 2, 3)와 수직 스티어러 3대(VCM1, 2, 3) 등 총 6대의 전자석이 s = 0.2m부터 8.0m 구간에 걸쳐 배치되어 있습니다. 빔 궤도 진단을 위한 센서로는 12대의 BPM(BPM 1~12)이 s = 1.0m부터 15.0m까지 전송선 전체에 걸쳐 균등하게 설치되어 있습니다. 스티어러 전자석은 전송선 전반부에 집중되어 있고, 진단 BPM은 끝단까지 배치되어 있음을 주목해 주십시오. 여러분의 제어기는 12개 BPM의 변위를 관측하여 6개 스티어러 전류를 정밀하게 인가해야 합니다.
\end{tcolorbox}
\vspace{0.2em}
\begin{tcolorbox}[colback=lightbg,colframe=bordergray,left=6pt,right=6pt,top=3pt,bottom=3pt,boxrule=0.6pt,arc=2.5pt]
\small
\noindent\textbf{\color{navy}교수학적 강조점 (Focus Point):}\enspace 제어기 관점에서 관측값은 12차원 BPM 벡터, 제어 행동은 6차원 전류 델타(ΔI) 벡터임을 명확히 연결 짓습니다.\par\vspace{0.3em}
\noindent\textbf{\color{successgreen}예상 질문 및 모범 답변 (Q\&A):}\enspace Q: 왜 스티어러는 6개인데 BPM은 12개인가요? A: 과결정(over-determined) 시스템이므로 완전 역행렬 대신 의사역행렬이나 최적화 기법이 필요합니다.\par\vspace{0.3em}
\noindent\textbf{\color{accent}다음 슬라이드 연결 멘트 (Transition):}\enspace 이 배치를 염두에 두고, 각 보정 전자석이 각 BPM 위치의 궤도를 얼마나 움직이는지, 즉 궤도 응답 행렬(ORM)을 유도해 보겠습니다.
\end{tcolorbox}
\clearpage
\hypertarget{slide.9}{}
\phantomsection
\addcontentsline{toc}{section}{슬라이드 9: 궤도 응답 행렬(Orbit Response Matrix, ORM) 유도}
\noindent\begin{tcolorbox}[colback=navy,colframe=navy,left=6pt,right=6pt,top=3pt,bottom=3pt,boxrule=0pt,arc=2pt]
{\large\bfseries\color{white}슬라이드 9:\enspace 궤도 응답 행렬(Orbit Response Matrix, ORM) 유도}\hfill{\footnotesize\bfseries\color{white!85!lightbg}Analytical Optics \enspace$|$\enspace 슬라이드 9 / 35}
\end{tcolorbox}
\vspace{0.2em}
\begin{tcolorbox}[colback=lightbg,colframe=bordergray,title={\footnotesize\bfseries\color{darkgray}슬라이드 시각 구성 요약},fontupper=\small,left=6pt,right=6pt,top=2.5pt,bottom=2.5pt,boxrule=0.6pt,arc=2.5pt]
단일 통과 빔라인 그린 함수, 감쇠 분모 (1 + 0.05·Δs)가 포함된 R\_ij 수식, 상류에서 하류로만 작용하는 인과율(BPM j가 보정자석 i보다 상류이면 R\_ij = 0).
\end{tcolorbox}
\vspace{0.2em}
\begin{tcolorbox}[colback=white,colframe=primaryblue,title={\footnotesize\bfseries\color{white}발표자 구두 대본 (Spoken Script)},fontupper=\small,left=6pt,right=6pt,top=3pt,bottom=3pt,boxrule=0.8pt,arc=2.5pt]
위치 s\_i에 위치한 보정용 스티어링 전자석이 각도 θ\_i만큼의 킥을 가하면, 입자 궤도에 Δx' = θ\_i의 기울기 불연속이 발생합니다. 하류의 s\_j 위치에 있는 빔 위치 모니터(BPM)에서 측정되는 궤도 변화는 빔라인 전달 그린 함수(Green's function)에 의해 결정됩니다. 단위 킥당 궤도 변화를 나타내는 궤도 응답 행렬(ORM) 원소 R\_ij는 R\_ij = √(β\_i·β\_j)·sin(ψ\_j - ψ\_i)로 정의됩니다. 우리 해커톤 전송선에서는 위상 전진율 2π×0.18 rad/m와 완만한 진폭 감쇠를 반영하여 R\_ij = 1.2·sin(2π·0.18·Δs) / (1 + 0.05·Δs)로 모델링되었습니다. 여기서 가장 중요한 물리적 법칙은 '인과율(causality)'입니다. 스티어러는 자신보다 하류(downstream, s\_j > s\_i)에 위치한 BPM에만 영향을 줄 수 있으며, 상류(upstream)에 있는 BPM에 대한 응답은 엄밀하게 0입니다. 이로 인해 6×12 크기의 ORM 행렬에는 계단 모양으로 0이 채워지며, 뒤쪽 보정자석일수록 영향을 주는 BPM 수가 줄어듭니다.
\end{tcolorbox}
\vspace{0.2em}
\begin{tcolorbox}[colback=lightbg,colframe=bordergray,left=6pt,right=6pt,top=3pt,bottom=3pt,boxrule=0.6pt,arc=2.5pt]
\small
\noindent\textbf{\color{navy}교수학적 강조점 (Focus Point):}\enspace 저장링의 폐궤도 공식(분모에 2·sin(πν) 포함)과 단일 통과 전송선의 인과적 그린 함수(sin 항) 간의 명확한 차이점을 환기합니다.\par\vspace{0.3em}
\noindent\textbf{\color{successgreen}예상 질문 및 모범 답변 (Q\&A):}\enspace Q: 상류 스티어러로 하류 오차를 잡을 수 있나요? A: 가능합니다. 반대로 하류 스티어러로는 상류의 왜곡을 물리적으로 되돌릴 수 없습니다.\par\vspace{0.3em}
\noindent\textbf{\color{accent}다음 슬라이드 연결 멘트 (Transition):}\enspace 실제 가속기 현장에서는 전원 공급 장치와 센서의 하드웨어적 한계가 존재합니다. 다음 슬라이드에서 이를 확인해보겠습니다.
\end{tcolorbox}
\clearpage
\hypertarget{slide.10}{}
\phantomsection
\addcontentsline{toc}{section}{슬라이드 10: 하드웨어 구동 한계 및 센서 노이즈 제약}
\noindent\begin{tcolorbox}[colback=navy,colframe=navy,left=6pt,right=6pt,top=3pt,bottom=3pt,boxrule=0pt,arc=2pt]
{\large\bfseries\color{white}슬라이드 10:\enspace 하드웨어 구동 한계 및 센서 노이즈 제약}\hfill{\footnotesize\bfseries\color{white!85!lightbg}Physical Limits \enspace$|$\enspace 슬라이드 10 / 35}
\end{tcolorbox}
\vspace{0.2em}
\begin{tcolorbox}[colback=lightbg,colframe=bordergray,title={\footnotesize\bfseries\color{darkgray}슬라이드 시각 구성 요약},fontupper=\small,left=6pt,right=6pt,top=2.5pt,bottom=2.5pt,boxrule=0.6pt,arc=2.5pt]
하드웨어 한계 요약: 총 코일 전류 허용치 [-5.0A, +5.0A], 단계당 슬루레이트 제한 [-0.8A, +0.8A], 센서 가우시안 노이즈 바닥 σ = 10 ㎛.
\end{tcolorbox}
\vspace{0.2em}
\begin{tcolorbox}[colback=white,colframe=primaryblue,title={\footnotesize\bfseries\color{white}발표자 구두 대본 (Spoken Script)},fontupper=\small,left=6pt,right=6pt,top=3pt,bottom=3pt,boxrule=0.8pt,arc=2.5pt]
이상적인 시뮬레이션에서 잘 작동하던 알고리즘이 실제 가속기에 적용되었을 때 실패하는 가장 큰 원인은 하드웨어 제약 조건을 간과하기 때문입니다. 본 벤치마크 환경에는 세 가지 엄격한 물리적 한계가 구현되어 있습니다. 첫째, '총 전류 한계'입니다. 전자석 코일의 정격 용량으로 인해 총 전류는 [-5.0 A, +5.0 A] 범위로 제한됩니다. 둘째, '슬루레이트(slew rate) 및 유도 기전력 한계'입니다. 코일 인덕턴스(L)와 전원 전압 한계(dI/dt = V/L)로 인해, 한 번의 제어 단계에서 인가할 수 있는 전류 증분은 최대 0.8 암페어로 엄격히 제한됩니다: ΔI ∈ [-0.8 A, +0.8 A]. 초기 궤도 왜곡을 보정하기 위해 필요한 누적 전류가 최대 2.84 A에 달하므로, 어떤 제어기도 1스텝 만에 수렴할 수는 없습니다. 이 스텝 제한 때문에 평가 시드 30개 중 10개는 최소 3스텝, 나머지 20개는 최소 4스텝이 필요하며 평균 하한은 3.7스텝입니다. 셋째, '센서 노이즈 바닥'입니다. 모든 BPM 측정값에는 표준편차 σ = 0.010 mm (10 ㎛)의 가우시안 지터가 포함됩니다. 10 ㎛ 이하의 잔여 오차를 억지로 잡으려 하면 노이즈를 증폭시키는 과적합이 발생합니다.
\end{tcolorbox}
\vspace{0.2em}
\begin{tcolorbox}[colback=lightbg,colframe=bordergray,left=6pt,right=6pt,top=3pt,bottom=3pt,boxrule=0.6pt,arc=2.5pt]
\small
\noindent\textbf{\color{navy}교수학적 강조점 (Focus Point):}\enspace 0.8A 스텝 제한으로 인해 최소 3~4스텝이 소요된다는 점을 강조하여, 해커톤 수렴 속도 점수의 물리적 한계를 인지시킵니다.\par\vspace{0.3em}
\noindent\textbf{\color{successgreen}예상 질문 및 모범 답변 (Q\&A):}\enspace Q: 제어기가 ΔI = 2.0 A를 출력하면 어떻게 되나요? A: 환경 내부에서 0.8 A로 강제 클리핑되지만, 바람직한 제어기는 자체적으로 경계를 준수해야 합니다.\par\vspace{0.3em}
\noindent\textbf{\color{accent}다음 슬라이드 연결 멘트 (Transition):}\enspace 하드웨어 한계에 더해, 전자석 철심 내부에서 일어나는 비선형 자기 포화 현상을 살펴보겠습니다.
\end{tcolorbox}
\clearpage
\hypertarget{slide.11}{}
\phantomsection
\addcontentsline{toc}{section}{슬라이드 11: 전자석 철심의 비선형 자기 포화(Magnetic Core Saturation)}
\noindent\begin{tcolorbox}[colback=navy,colframe=navy,left=6pt,right=6pt,top=3pt,bottom=3pt,boxrule=0pt,arc=2pt]
{\large\bfseries\color{white}슬라이드 11:\enspace 전자석 철심의 비선형 자기 포화(Magnetic Core Saturation)}\hfill{\footnotesize\bfseries\color{white!85!lightbg}Non-Linear Physics \enspace$|$\enspace 슬라이드 11 / 35}
\end{tcolorbox}
\vspace{0.2em}
\begin{tcolorbox}[colback=lightbg,colframe=bordergray,title={\footnotesize\bfseries\color{darkgray}슬라이드 시각 구성 요약},fontupper=\small,left=6pt,right=6pt,top=2.5pt,bottom=2.5pt,boxrule=0.6pt,arc=2.5pt]
B-H 자기 포화 곡선, 3차 비선형 수식 Y = I·R - 0.015·(I$^3$·R) + Y$_0$ + ε, 선형 응답과 비선형 포화 응답 곡선의 편차 비교.
\end{tcolorbox}
\vspace{0.2em}
\begin{tcolorbox}[colback=white,colframe=primaryblue,title={\footnotesize\bfseries\color{white}발표자 구두 대본 (Spoken Script)},fontupper=\small,left=6pt,right=6pt,top=3pt,bottom=3pt,boxrule=0.8pt,arc=2.5pt]
기초 가속기 물리학에서는 암페어 법칙에 따라 전자석 자장 B가 인가 전류 I에 정비례한다고 가정합니다: B = $\mu_0$·N·I / g. 하지만 실제 전자석 철심은 강자성체인 규소강판 등으로 제작됩니다. 전류가 낮을 때는 자구(magnetic domain)가 쉽게 정렬되지만, 전류가 커질수록 투자율이 떨어져 자장이 전류만큼 증가하지 못합니다. 우리 환경에서는 이를 3차 소프트 포화 항으로 정밀하게 모사했습니다: Y = I·R - 0.015·(I$^3$·R) + Y$_0$ + ε. 이 수식이 의미하는 바를 주목해 주십시오. 손실 비율은 0.015·I$^2$입니다. 2 A에서 6\%, 우리 보정에 필요한 최대 전류인 2.84 A에서 12\%, 4 A에서는 24\%입니다. 선형 광학 모델(ORM)에만 의존하는 제어기는 필요한 킥의 크기를 과소평가하여 보정이 부족해지고, 그만큼 스텝이 더 필요합니다. 바로 이 비선형 포화 효과를 사전에 학습하고 보상하기 위해 머신러닝 대리 모델이 필요한 것입니다.
\end{tcolorbox}
\vspace{0.2em}
\begin{tcolorbox}[colback=lightbg,colframe=bordergray,left=6pt,right=6pt,top=3pt,bottom=3pt,boxrule=0.6pt,arc=2.5pt]
\small
\noindent\textbf{\color{navy}교수학적 강조점 (Focus Point):}\enspace 초기 궤도 보정에 최대 2.84 A의 전류가 필요하므로, 3차 포화 효과가 극대화되는 비선형 영역에서 제어가 이루어짐을 역설합니다.\par\vspace{0.3em}
\noindent\textbf{\color{successgreen}예상 질문 및 모범 답변 (Q\&A):}\enspace Q: 고전 SVD 제어기는 이 비선형 포화 상태에서 수렴하지 못하나요? A: 게인(α=0.45)을 낮추면 수렴은 하지만, 9.4스텝이나 걸려 고득점을 얻기 어렵습니다.\par\vspace{0.3em}
\noindent\textbf{\color{accent}다음 슬라이드 연결 멘트 (Transition):}\enspace 이제 머신러닝의 기초로 넘어가서, 인공신경망이 어떻게 이러한 복잡한 비선형 함수를 학습하는지 원리를 살펴보겠습니다.
\end{tcolorbox}
\clearpage
\hypertarget{slide.12}{}
\phantomsection
\addcontentsline{toc}{section}{슬라이드 12: 인공 뉴런(Artificial Neuron)과 신경망 계층 구성}
\noindent\begin{tcolorbox}[colback=navy,colframe=navy,left=6pt,right=6pt,top=3pt,bottom=3pt,boxrule=0pt,arc=2pt]
{\large\bfseries\color{white}슬라이드 12:\enspace 인공 뉴런(Artificial Neuron)과 신경망 계층 구성}\hfill{\footnotesize\bfseries\color{white!85!lightbg}ML Foundations \enspace$|$\enspace 슬라이드 12 / 35}
\end{tcolorbox}
\vspace{0.2em}
\begin{tcolorbox}[colback=lightbg,colframe=bordergray,title={\footnotesize\bfseries\color{darkgray}슬라이드 시각 구성 요약},fontupper=\small,left=6pt,right=6pt,top=2.5pt,bottom=2.5pt,boxrule=0.6pt,arc=2.5pt]
인공 뉴런 구조: 입력 x\_i, 시냅스 가중치 w\_i, 편향 b, 선형 결합 z = w$^T$x + b, 활성화 함수 σ(z), 다층 퍼셉트론(MLP) 레이어 적층 모식도.
\end{tcolorbox}
\vspace{0.2em}
\begin{tcolorbox}[colback=white,colframe=primaryblue,title={\footnotesize\bfseries\color{white}발표자 구두 대본 (Spoken Script)},fontupper=\small,left=6pt,right=6pt,top=3pt,bottom=3pt,boxrule=0.8pt,arc=2.5pt]
딥러닝의 가장 기본 단위는 1943년 매컬럭과 피츠가 제안한 '인공 뉴런(artificial neuron)'입니다. 뉴런은 n차원 입력 벡터 x를 받습니다. 각 입력값은 학습 가능한 시냅스 가중치 w\_i와 곱해진 후 합산되며, 임계값을 조절하는 편향 b가 더해집니다: z = w$^T$x + b. 이 선형 결합 값 z는 비선형 활성화 함수 σ(z)를 통과하여 최종 출력 a를 생성합니다. 이러한 뉴런들을 병렬로 배치하면 밀집 레이어(dense layer) h = σ(W x + b)가 형성되며, 여러 개의 레이어를 순차적으로 적층하면 '다층 퍼셉트론(Multi-Layer Perceptron, MLP)'이 완성됩니다. 시벤코(Cybenko)와 호닉(Hornik)의 보편 근사 정리(Universal Approximation Theorem)에 따르면, 비선형 활성화 함수를 가진 다층 신경망은 적절한 가중치만 주어진다면 어떠한 연속 비선형 물리 함수도 임의의 정밀도로 근사할 수 있습니다.
\end{tcolorbox}
\vspace{0.2em}
\begin{tcolorbox}[colback=lightbg,colframe=bordergray,left=6pt,right=6pt,top=3pt,bottom=3pt,boxrule=0.6pt,arc=2.5pt]
\small
\noindent\textbf{\color{navy}교수학적 강조점 (Focus Point):}\enspace 선형 변환(W x + b)과 비선형 활성화 함수 σ(z)의 명확한 역할 분담을 설명합니다.\par\vspace{0.3em}
\noindent\textbf{\color{successgreen}예상 질문 및 모범 답변 (Q\&A):}\enspace Q: 활성화 함수 σ가 선형이면 어떻게 되나요? A: 아무리 많은 레이어를 쌓아도 단 하나의 행렬 곱으로 축퇴되어 비선형성을 전혀 표현할 수 없습니다.\par\vspace{0.3em}
\noindent\textbf{\color{accent}다음 슬라이드 연결 멘트 (Transition):}\enspace 그렇다면 활성화 함수는 신경망 내부에서 구체적으로 어떤 역할을 하며, 어떤 함수들이 사용되는지 확인해보겠습니다.
\end{tcolorbox}
\clearpage
\hypertarget{slide.13}{}
\phantomsection
\addcontentsline{toc}{section}{슬라이드 13: 신경망 활성화 함수(Activation Functions)의 역할과 분류}
\noindent\begin{tcolorbox}[colback=navy,colframe=navy,left=6pt,right=6pt,top=3pt,bottom=3pt,boxrule=0pt,arc=2pt]
{\large\bfseries\color{white}슬라이드 13:\enspace 신경망 활성화 함수(Activation Functions)의 역할과 분류}\hfill{\footnotesize\bfseries\color{white!85!lightbg}ML Foundations \enspace$|$\enspace 슬라이드 13 / 35}
\end{tcolorbox}
\vspace{0.2em}
\begin{tcolorbox}[colback=lightbg,colframe=bordergray,title={\footnotesize\bfseries\color{darkgray}슬라이드 시각 구성 요약},fontupper=\small,left=6pt,right=6pt,top=2.5pt,bottom=2.5pt,boxrule=0.6pt,arc=2.5pt]
활성화 함수의 4대 핵심 역할, 주요 활성화 함수 수식 및 특성(ReLU, Leaky ReLU, Tanh, Sigmoid, GELU, Linear), 함수값 및 도함수 비교 그래프.
\end{tcolorbox}
\vspace{0.2em}
\begin{tcolorbox}[colback=white,colframe=primaryblue,title={\footnotesize\bfseries\color{white}발표자 구두 대본 (Spoken Script)},fontupper=\small,left=6pt,right=6pt,top=3pt,bottom=3pt,boxrule=0.8pt,arc=2.5pt]
가속기 신경망에서 활성화 함수는 네 가지 본질적인 역할을 수행합니다. 첫째, '선형성 타파'입니다. 활성화 함수가 없다면 아무리 깊은 신경망도 하나의 선형 행렬로 붕괴됩니다. 둘째, '비선형 물리 동역학의 근사'입니다. 자기 포화나 6극자 왜곡을 정확히 표현할 수 있게 합니다. 셋째, '역전파 그래디언트 제어'입니다. 도함수 σ'(z)는 오차 신호가 상류 레이어로 전달되는 게이트 역할을 하며, 기울기 소실(vanishing gradient)을 방지합니다. 넷째, '출력 공간의 물리적 연속성 보장'입니다. 은닉층에는 비선형성이 필수적이지만, 빔 궤도 변위(mm)를 예측하는 출력층은 반드시 선형 항등 함수(linear identity)를 사용하여 인위적인 값 클리핑을 막아야 합니다. 가장 널리 쓰이는 ReLU는 z > 0에서 도함수가 항상 1이므로 기울기 소실을 완벽히 방지하며 계산 효율이 매우 뛰어납니다.
\end{tcolorbox}
\vspace{0.2em}
\begin{tcolorbox}[colback=lightbg,colframe=bordergray,left=6pt,right=6pt,top=3pt,bottom=3pt,boxrule=0.6pt,arc=2.5pt]
\small
\noindent\textbf{\color{navy}교수학적 강조점 (Focus Point):}\enspace 네 번째 역할인 '출력층의 선형 활성화 함수 필수 사용'을 거듭 강조합니다. 출력에 Sigmoid 등을 적용하면 빔 변위가 인위적으로 잘리게 됩니다.\par\vspace{0.3em}
\noindent\textbf{\color{successgreen}예상 질문 및 모범 답변 (Q\&A):}\enspace Q: 'Dying ReLU' 현상이란 무엇인가요? A: 입력값이 0 이하가 되면 기울기가 0이 되어 가중치가 더 이상 업데이트되지 않는 현상이며, Leaky ReLU로 완화할 수 있습니다.\par\vspace{0.3em}
\noindent\textbf{\color{accent}다음 슬라이드 연결 멘트 (Transition):}\enspace 이러한 활성화 함수의 거동을 입자 전달 행렬을 모사하는 구체적인 예제를 통해 시각적으로 확인해보겠습니다.
\end{tcolorbox}
\clearpage
\hypertarget{slide.14}{}
\phantomsection
\addcontentsline{toc}{section}{슬라이드 14: 신경망 대리 모델 예제: 입자 전달 행렬(Transfer Matrix M) 학습}
\noindent\begin{tcolorbox}[colback=navy,colframe=navy,left=6pt,right=6pt,top=3pt,bottom=3pt,boxrule=0pt,arc=2pt]
{\large\bfseries\color{white}슬라이드 14:\enspace 신경망 대리 모델 예제: 입자 전달 행렬(Transfer Matrix M) 학습}\hfill{\footnotesize\bfseries\color{white!85!lightbg}ML Foundations \enspace$|$\enspace 슬라이드 14 / 35}
\end{tcolorbox}
\vspace{0.2em}
\begin{tcolorbox}[colback=lightbg,colframe=bordergray,title={\footnotesize\bfseries\color{darkgray}슬라이드 시각 구성 요약},fontupper=\small,left=6pt,right=6pt,top=2.5pt,bottom=2.5pt,boxrule=0.6pt,arc=2.5pt]
초기 위상공간 (x$_0$, x'$_0$)에서 최종 (x\_f, x'\_f)로의 매핑 수식, 벡터 표현 $\hat{z}_f$ = f\_θ(z$_0$), 심플렉틱 조건 $\det(M)$ = 1, 확대된 TikZ 신경망 구조도.
\end{tcolorbox}
\vspace{0.2em}
\begin{tcolorbox}[colback=white,colframe=primaryblue,title={\footnotesize\bfseries\color{white}발표자 구두 대본 (Spoken Script)},fontupper=\small,left=6pt,right=6pt,top=3pt,bottom=3pt,boxrule=0.8pt,arc=2.5pt]
머신러닝과 빔 광학을 직접 연결해보기 위해, 선형 전달 행렬 M을 신경망으로 대리 모델링하는 예제를 구성했습니다. 선형 빔 광학에서 초기 위치 s$_0$의 횡방향 좌표와 각도 (x$_0$, x'$_0$)는 2×2 전달 행렬 M을 통해 최종 위치 (x\_f, x'\_f)로 변환됩니다: [x\_f, x'\_f]$^T$ = M · [x$_0$, x'$_0$]$^T$. 이때 리우빌 정리(Liouville's theorem)에 의해 행렬식 $\det M$ = 1을 만족하며 위상 공간 면적이 보존됩니다. 우리는 해석적 행렬 곱 대신 2차원에서 2차원으로 매핑하는 신경망 f\_θ를 구성했습니다: $\hat{z}_f$ = W$^{(3)}$ σ( W$^{(2)}$ σ( W$^{(1)}$ z$_0$ + b$^{(1)}$ ) + b$^{(2)}$ ) + b$^{(3)}$. 수천 개의 트래킹 입자 쌍에 대해 평균제곱오차(MSE) 손실 함수로 이 네트워크를 학습시킵니다. 오른쪽 구조도를 보시면 2개의 입력 뉴런, 4개의 뉴런을 갖는 2개의 은닉층, 그리고 최종 좌표를 예측하는 2개의 출력 뉴런으로 이루어져 있습니다. 하단에 강조되어 있듯이, 최종 출력 뉴런의 활성화 함수는 반드시 선형 항등 함수여야 합니다.
\end{tcolorbox}
\vspace{0.2em}
\begin{tcolorbox}[colback=lightbg,colframe=bordergray,left=6pt,right=6pt,top=3pt,bottom=3pt,boxrule=0.6pt,arc=2.5pt]
\small
\noindent\textbf{\color{navy}교수학적 강조점 (Focus Point):}\enspace 컴팩트한 벡터 표기법 $\hat{z}_f$ = f\_θ(z$_0$)를 설명하고, 신경망 야코비 행렬식 $\det(J)$가 위상공간 면적 보존을 나타냄을 언급합니다.\par\vspace{0.3em}
\noindent\textbf{\color{successgreen}예상 질문 및 모범 답변 (Q\&A):}\enspace Q: 일반 MLP가 $\det(M)$ = 1을 완벽히 보존할 수 있나요? A: 일반 MLP는 오차 범위 내에서 근사할 뿐이며, 엄밀한 보존을 위해서는 해밀토니안 신경망(HNN)이나 PINN 구조가 필요합니다.\par\vspace{0.3em}
\noindent\textbf{\color{accent}다음 슬라이드 연결 멘트 (Transition):}\enspace 이 전달 행렬 대리 모델의 은닉층에 서로 다른 활성화 함수를 적용하면 위상 공간에서 어떤 기하학적 왜곡이 일어나는지 살펴보겠습니다.
\end{tcolorbox}
\clearpage
\hypertarget{slide.15}{}
\phantomsection
\addcontentsline{toc}{section}{슬라이드 15: 빔 전송 대리 모델에서 활성화 함수별 위상 공간 변형}
\noindent\begin{tcolorbox}[colback=navy,colframe=navy,left=6pt,right=6pt,top=3pt,bottom=3pt,boxrule=0pt,arc=2pt]
{\large\bfseries\color{white}슬라이드 15:\enspace 빔 전송 대리 모델에서 활성화 함수별 위상 공간 변형}\hfill{\footnotesize\bfseries\color{white!85!lightbg}ML Foundations \enspace$|$\enspace 슬라이드 15 / 35}
\end{tcolorbox}
\vspace{0.2em}
\begin{tcolorbox}[colback=lightbg,colframe=bordergray,title={\footnotesize\bfseries\color{darkgray}슬라이드 시각 구성 요약},fontupper=\small,left=6pt,right=6pt,top=2.5pt,bottom=2.5pt,boxrule=0.6pt,arc=2.5pt]
선형(Linear), ReLU, Tanh 활성화 함수가 위상공간 타원에 미치는 기하학적 영향 비교: 행렬 축퇴, 다면체 꺾임, 에미턴스 수축 현상.
\end{tcolorbox}
\vspace{0.2em}
\begin{tcolorbox}[colback=white,colframe=primaryblue,title={\footnotesize\bfseries\color{white}발표자 구두 대본 (Spoken Script)},fontupper=\small,left=6pt,right=6pt,top=3pt,bottom=3pt,boxrule=0.8pt,arc=2.5pt]
이 슬라이드는 전달 행렬을 학습할 때 활성화 함수 선택이 위상 공간에 미치는 심오한 기하학적 영향을 보여줍니다. 첫째, '선형 활성화 함수'를 은닉층에 사용하면 모든 계층이 단 하나의 유효 행렬 곱 W\_eff x + b\_net으로 축퇴됩니다. 수학적으로 선형 전달 행렬과 완전히 동일하므로 매끄러운 빔 타원과 심플렉틱 면적 보존을 완벽히 유지하지만, 3차 자기 포화나 6극자 왜곡을 전혀 표현할 수 없습니다. 둘째, 'ReLU'는 조각별 선형(piecewise linear) 함수이므로 위상 공간을 초평면으로 분할합니다. 그 결과 매끄러운 빔 타원이 모서리가 각진 '다면체 다각형(faceted polygon)'으로 변형됩니다! 더욱이 뉴런이 비활성화(z < 0)되는 음의 영역에서는 야코비 행렬식이 0으로 붕괴되어 리우빌 정리를 심각하게 위반합니다. 셋째, 'Tanh'는 원점 부근에서는 tanh(z) ≈ z로 선형 광학을 잘 모사하지만, 진폭이 큰 빔 헤일로(halo) 영역에서는 ±1로 포화되어 타원을 S자 형태로 찌그러뜨리고 빔 에미턴스를 인위적으로 축소시키는 왜곡을 일으킵니다.
\end{tcolorbox}
\vspace{0.2em}
\begin{tcolorbox}[colback=lightbg,colframe=bordergray,left=6pt,right=6pt,top=3pt,bottom=3pt,boxrule=0.6pt,arc=2.5pt]
\small
\noindent\textbf{\color{navy}교수학적 강조점 (Focus Point):}\enspace 선형은 비선형을 못 배우고, ReLU는 타원을 각지게 만들며, Tanh는 에미턴스를 왜곡시킨다는 물리적 직관을 명확히 대조합니다. 컨테이너에 포함된 보조 노트북 Transfer\_Matrix\_Surrogate\_Phase\_Space.ipynb에서 각 효과를 위상공간 그림으로 확인할 수 있음을 안내합니다.\par\vspace{0.3em}
\noindent\textbf{\color{successgreen}예상 질문 및 모범 답변 (Q\&A):}\enspace Q: 이번 해커톤 MLP에는 어떤 활성화 함수가 가장 적합한가요? A: 기본값인 relu에서 시작하되 tanh를 꼭 시도해 보십시오. 이 응답은 매끄러운 함수이므로, lbfgs로 학습한 작은 tanh 신경망은 0.012 mm에 도달하는 반면 기본 설정은 0.109 mm에서 멈춥니다. MLPRegressor의 출력층은 항상 선형입니다.\par\vspace{0.3em}
\noindent\textbf{\color{accent}다음 슬라이드 연결 멘트 (Transition):}\enspace 신경망 구조를 정의한 후, 수만 개의 파라미터를 어떻게 최적화하는지 최신 최적화 알고리즘들을 살펴보겠습니다.
\end{tcolorbox}
\clearpage
\hypertarget{slide.16}{}
\phantomsection
\addcontentsline{toc}{section}{슬라이드 16: 머신러닝 최적화 알고리즘: 경사하강법부터 Adam 및 L-BFGS-B까지}
\noindent\begin{tcolorbox}[colback=navy,colframe=navy,left=6pt,right=6pt,top=3pt,bottom=3pt,boxrule=0pt,arc=2pt]
{\large\bfseries\color{white}슬라이드 16:\enspace 머신러닝 최적화 알고리즘: 경사하강법부터 Adam 및 L-BFGS-B까지}\hfill{\footnotesize\bfseries\color{white!85!lightbg}ML Foundations \enspace$|$\enspace 슬라이드 16 / 35}
\end{tcolorbox}
\vspace{0.2em}
\begin{tcolorbox}[colback=lightbg,colframe=bordergray,title={\footnotesize\bfseries\color{darkgray}슬라이드 시각 구성 요약},fontupper=\small,left=6pt,right=6pt,top=2.5pt,bottom=2.5pt,boxrule=0.6pt,arc=2.5pt]
1차 경사하강법 및 준뉴턴 최적화 비교 카드: SGD, 모멘텀, Adam(1차 및 2차 모멘트 적응형 스케일링), 박스 경계 준뉴턴 L-BFGS-B.
\end{tcolorbox}
\vspace{0.2em}
\begin{tcolorbox}[colback=white,colframe=primaryblue,title={\footnotesize\bfseries\color{white}발표자 구두 대본 (Spoken Script)},fontupper=\small,left=6pt,right=6pt,top=3pt,bottom=3pt,boxrule=0.8pt,arc=2.5pt]
신경망 학습이란 손실 함수 L(θ)를 최소화하는 파라미터 θ를 찾아가는 과정입니다. 이 슬라이드에서는 4가지 주요 최적화 알고리즘을 비교합니다. 첫째는 '확률적 경사하강법(SGD)'입니다: 최급강하 방향으로 θ ← θ - η∇L 만큼 이동합니다. 단순하지만 곡률이 큰 협곡 지형에서 발진하며 수렴이 느립니다. 둘째는 '모멘텀(Momentum) SGD'입니다: 이전 단계의 관성 벡터 γ·v를 더해주어 진동을 억제하고 일관된 경사 방향으로 빠르게 가속합니다. 셋째는 'Adam(Adaptive Moment Estimation)'입니다: 1차 모멘트(기울기 평균)와 2차 모멘트(기울기 제곱 평균)를 동시에 추적하여, 파라미터별로 학습률을 적응형으로 스케일링합니다. 민감한 파라미터의 스텝은 줄이고 둔감한 파라미터의 스텝은 키워주므로, 신경망 학습의 일반적인 기본 선택으로 사용됩니다. 넷째는 'L-BFGS-B'입니다: 헤시안 역행렬을 근사하는 2차 준뉴턴(quasi-Newton) 기법입니다. 오프라인 학습에 쓰이는 Adam과 달리, L-BFGS-B는 해커톤 제어기에서 매 스텝 대리 모델을 역산(inversion)하여 다음 킥을 도출할 때 사용됩니다.
\end{tcolorbox}
\vspace{0.2em}
\begin{tcolorbox}[colback=lightbg,colframe=bordergray,left=6pt,right=6pt,top=3pt,bottom=3pt,boxrule=0.6pt,arc=2.5pt]
\small
\noindent\textbf{\color{navy}교수학적 강조점 (Focus Point):}\enspace 역할 구분을 명확히 합니다: 기준 노트북에서는 Adam이 오프라인에서 대리 모델을 '학습'시키고, L-BFGS-B가 온라인에서 대리 모델을 '역산'하여 제어 킥을 찾습니다. L-BFGS 계열로 작은 신경망을 학습시킬 수도 있음을 언급합니다.\par\vspace{0.3em}
\noindent\textbf{\color{successgreen}예상 질문 및 모범 답변 (Q\&A):}\enspace Q: MLP 학습 자체에 L-BFGS를 쓰면 안 되나요? A: 쓸 수 있고, 이번 과제에서는 꼭 시도해 볼 만합니다. 매끄러운 학습 샷이 1,200개뿐이므로 전체 배치 방식인 solver='lbfgs'는 노이즈 바닥까지 도달하는 반면, 기본 Adam 학습은 일찍 멈춥니다. 미니배치 노이즈가 있는 대규모 데이터에서는 Adam이 더 적합합니다.\par\vspace{0.3em}
\noindent\textbf{\color{accent}다음 슬라이드 연결 멘트 (Transition):}\enspace 그렇다면 가속기에서는 MLP만 사용될까요? 전혀 아닙니다. 현대 가속기 물리에 특화된 고급 신경망 구조들을 살펴보겠습니다.
\end{tcolorbox}
\clearpage
\hypertarget{slide.17}{}
\phantomsection
\addcontentsline{toc}{section}{슬라이드 17: MLP를 넘어선 가속기 특화 신경망 구조: PINN, PRNN, VAE, U-Net}
\noindent\begin{tcolorbox}[colback=navy,colframe=navy,left=6pt,right=6pt,top=3pt,bottom=3pt,boxrule=0pt,arc=2pt]
{\large\bfseries\color{white}슬라이드 17:\enspace MLP를 넘어선 가속기 특화 신경망 구조: PINN, PRNN, VAE, U-Net}\hfill{\footnotesize\bfseries\color{white!85!lightbg}ML Foundations \enspace$|$\enspace 슬라이드 17 / 35}
\end{tcolorbox}
\vspace{0.2em}
\begin{tcolorbox}[colback=lightbg,colframe=bordergray,title={\footnotesize\bfseries\color{darkgray}슬라이드 시각 구성 요약},fontupper=\small,left=6pt,right=6pt,top=2.5pt,bottom=2.5pt,boxrule=0.6pt,arc=2.5pt]
4분면 패널 레이아웃: PINN(물리 미분방정식 손실 통합), PRNN/RNN(시간적 동역학 및 히스테리시스 메모리), VAE(확률적 잠재 공간 및 이상 진단), U-Net(다중 스케일 스킵 연결 및 빔 영상 복원).
\end{tcolorbox}
\vspace{0.2em}
\begin{tcolorbox}[colback=white,colframe=primaryblue,title={\footnotesize\bfseries\color{white}발표자 구두 대본 (Spoken Script)},fontupper=\small,left=6pt,right=6pt,top=3pt,bottom=3pt,boxrule=0.8pt,arc=2.5pt]
단순한 MLP는 정적 회귀 문제에는 우수하지만, 실제 가속기 현장의 복잡한 물리 법칙을 만족시키기에는 한계가 있습니다. 이 슬라이드에서는 가속기 물리에 최적화된 4가지 고급 딥러닝 아키텍처를 소개합니다. 첫째는 '물리 정보 신경망(PINN)'입니다: 힐 방정식이나 포아송 공간전하 방정식 같은 지배 미분방정식을 자동미분(autograd)을 통해 손실 함수에 직접 통합합니다: L\_PINN = L\_data + λ·||D[f](x)||$^2$. 이를 통해 데이터가 부족한 드리프트 구간에서도 에너지 및 심플렉틱 보존 법칙을 완벽히 만족합니다. 둘째는 '물리 유도 순환신경망(PRNN/RNN)'입니다: 내부 은닉 상태 h\_t를 통해 이전 시점의 이력을 기억합니다. 정적 입력만 처리하는 MLP와 달리, 전자석 충방전 시 발생하는 비선형 자기 히스테리시스 루프, 와전류(eddy current) 지연, 초전도 공동의 감쇠 현상을 정확히 추적합니다. 셋째는 '변이형 오토인코더(VAE)'입니다: 입력 데이터를 연속적인 가우시안 잠재 공간 z ~ N(μ, Σ)으로 압축하고 ELBO를 최대화합니다. 6차원 입자 다발 위상 공간 압축 및 빔 정지 전조 증상을 비지도 학습으로 감지하는 데 탁월합니다. 넷째는 'U-Net 및 합성곱 신경망(CNN)'입니다: 공간적 병진 대칭성을 활용하며, 축소 경로와 확장 경로를 연결하는 스킵 연결(skip connection)을 통해 YAG/OTR 형광 스크린 빔 프로파일 영상 복원, 빔 헤일로 분할, 3차원 자장 맵 보간을 서브픽셀 정밀도로 수행합니다.
\end{tcolorbox}
\vspace{0.2em}
\begin{tcolorbox}[colback=lightbg,colframe=bordergray,left=6pt,right=6pt,top=3pt,bottom=3pt,boxrule=0.6pt,arc=2.5pt]
\small
\noindent\textbf{\color{navy}교수학적 강조점 (Focus Point):}\enspace 각 사분면을 시계 방향으로 짚어가며, 각 구조가 단순 MLP의 어떤 물리적 한계점을 극복했는지 명확히 설명합니다.\par\vspace{0.3em}
\noindent\textbf{\color{successgreen}예상 질문 및 모범 답변 (Q\&A):}\enspace Q: 폐루프 궤도 제어에 PINN을 대리 모델로 쓸 수 있나요? A: 가능합니다. 보정 킥이 매우 커서 훈련 데이터 범위를 벗어날 때 PINN이 훨씬 물리적으로 신뢰할 수 있는 예측을 제공합니다.\par\vspace{0.3em}
\noindent\textbf{\color{accent}다음 슬라이드 연결 멘트 (Transition):}\enspace 이제 오늘 미니 해커톤의 첫 번째 핵심 트랙인 지도학습 대리 모델 역산 제어로 들어가 보겠습니다.
\end{tcolorbox}
\clearpage
\hypertarget{slide.18}{}
\phantomsection
\addcontentsline{toc}{section}{슬라이드 18: 지도학습: 디지털 트윈 대리 모델링 및 하이퍼파라미터}
\noindent\begin{tcolorbox}[colback=navy,colframe=navy,left=6pt,right=6pt,top=3pt,bottom=3pt,boxrule=0pt,arc=2pt]
{\large\bfseries\color{white}슬라이드 18:\enspace 지도학습: 디지털 트윈 대리 모델링 및 하이퍼파라미터}\hfill{\footnotesize\bfseries\color{white!85!lightbg}Paradigm 1: Supervised Learning \enspace$|$\enspace 슬라이드 18 / 35}
\end{tcolorbox}
\vspace{0.2em}
\begin{tcolorbox}[colback=lightbg,colframe=bordergray,title={\footnotesize\bfseries\color{darkgray}슬라이드 시각 구성 요약},fontupper=\small,left=6pt,right=6pt,top=2.5pt,bottom=2.5pt,boxrule=0.6pt,arc=2.5pt]
좌측: 핵심 용어(가중치, 손실 함수, 역전파, 에포크, 배치, 조기 종료). 우측: 참가자 튜닝 가능 변수(TODO 1)가 빨간색으로 강조된 하이퍼파라미터 분류표.
\end{tcolorbox}
\vspace{0.2em}
\begin{tcolorbox}[colback=white,colframe=primaryblue,title={\footnotesize\bfseries\color{white}발표자 구두 대본 (Spoken Script)},fontupper=\small,left=6pt,right=6pt,top=3pt,bottom=3pt,boxrule=0.8pt,arc=2.5pt]
첫 번째 패러다임은 가속기 디지털 트윈의 표준 기법인 '지도학습 순방향 대리 모델링'입니다. 좌측에는 모델 학습의 핵심 메커니즘을 정리했습니다. 가중치 W와 편향 b는 손실 함수의 기울기를 역전파하여 반복 수정됩니다. 에포크(epoch)는 1,200개의 훈련 데이터 전체를 1회 완독하는 주기이며, 배치(batch)는 1회 가중치 업데이트에 사용되는 샘플 묶음입니다. 조기 종료(early stopping)는 검증 손실이 더 이상 줄어들지 않을 때 학습을 멈춰 센서 노이즈 과적합을 방지합니다. 우측 표에는 여러분이 노트북 TODO 1에서 직접 수정할 수 있는 핵심 하이퍼파라미터들이 빨간색으로 강조되어 있습니다. 학습률(기본 10$^-$$^3$), 배치 크기(32~128), 은닉층 구조(기본 128, 64, 32), 활성화 함수(relu/tanh), 솔버(adam/lbfgs), L2 규제 계수(10$^-$⁴), 최대 에포크 수 등이 포함됩니다. 이 하이퍼파라미터들을 어떻게 최적화하느냐에 따라 여러분의 디지털 트윈 정밀도가 결정됩니다.
\end{tcolorbox}
\vspace{0.2em}
\begin{tcolorbox}[colback=lightbg,colframe=bordergray,left=6pt,right=6pt,top=3pt,bottom=3pt,boxrule=0.6pt,arc=2.5pt]
\small
\noindent\textbf{\color{navy}교수학적 강조점 (Focus Point):}\enspace 빨간색 별표(*)로 표시된 행들이 참가자 노트북 TODO 1에서 직접 수정할 수 있는 변수임을 강조합니다.\par\vspace{0.3em}
\noindent\textbf{\color{successgreen}예상 질문 및 모범 답변 (Q\&A):}\enspace Q: 배치 크기가 너무 작으면 어떻게 되나요? A: 그래디언트 노이즈가 커지고 한 에포크에 필요한 업데이트 횟수가 늘어납니다. 배치가 아주 크면 이러한 규제 효과가 사라지며, 전체 배치 학습이 바로 solver='lbfgs'의 방식입니다.\par\vspace{0.3em}
\noindent\textbf{\color{accent}다음 슬라이드 연결 멘트 (Transition):}\enspace 그렇다면 이 디지털 트윈 신경망의 입출력 구조와 계층 배치는 어떻게 이루어져 있는지 확인해보겠습니다.
\end{tcolorbox}
\clearpage
\hypertarget{slide.19}{}
\phantomsection
\addcontentsline{toc}{section}{슬라이드 19: 신경망 대리 모델 디지털 트윈 구조}
\noindent\begin{tcolorbox}[colback=navy,colframe=navy,left=6pt,right=6pt,top=3pt,bottom=3pt,boxrule=0pt,arc=2pt]
{\large\bfseries\color{white}슬라이드 19:\enspace 신경망 대리 모델 디지털 트윈 구조}\hfill{\footnotesize\bfseries\color{white!85!lightbg}Paradigm 1: Supervised Learning \enspace$|$\enspace 슬라이드 19 / 35}
\end{tcolorbox}
\vspace{0.2em}
\begin{tcolorbox}[colback=lightbg,colframe=bordergray,title={\footnotesize\bfseries\color{darkgray}슬라이드 시각 구성 요약},fontupper=\small,left=6pt,right=6pt,top=2.5pt,bottom=2.5pt,boxrule=0.6pt,arc=2.5pt]
신경망 디지털 트윈 구조도: 6개 입력 전류 [I$_1$~I$_6$], 3개 은닉층(128, 64, 32 뉴런)과 ReLU 활성화, 12개 출력 예측 궤도 [$\hat{Y}$$_1$~$\hat{Y}$$_1$$_2$].
\end{tcolorbox}
\vspace{0.2em}
\begin{tcolorbox}[colback=white,colframe=primaryblue,title={\footnotesize\bfseries\color{white}발표자 구두 대본 (Spoken Script)},fontupper=\small,left=6pt,right=6pt,top=3pt,bottom=3pt,boxrule=0.8pt,arc=2.5pt]
이 슬라이드는 해커톤에서 사용되는 신경망 대리 모델의 구체적인 구조도입니다. 6개의 보정 전자석 전류를 입력받아 12개 BPM의 궤도 변위를 실시간 예측하는 R$^6$ → R¹$^2$ 매핑 구조입니다. 입력층은 6개의 수평/수직 스티어러 전류 [I$_1$, ..., I$_6$]를 받습니다. 이 신호는 세 개의 밀집 은닉층을 거칩니다: 첫 번째 은닉층은 6개의 입력을 128개의 고차원 특징 공간으로 확장하고, 두 번째 층은 64개로 압축하며, 세 번째 층은 32개로 정제합니다. 세 은닉층 모두 ReLU 활성화 함수를 사용합니다. 마지막 출력층은 32개의 추상 특징을 12개 BPM 예측 변위값 [$\hat{Y}$$_1$, ..., $\hat{Y}$$_1$$_2$]으로 선형 투영합니다. 이 네트워크의 총 학습 파라미터는 11,628개입니다. 선형 ORM 행렬이 단 72개의 파라미터만 갖는 것에 비해, 11,628개의 가중치는 3차 비선형 자기 포화 효과를 표현할 수 있는 용량을 제공합니다. 다만 다음 슬라이드에서 보듯이 용량만으로는 충분하지 않습니다.
\end{tcolorbox}
\vspace{0.2em}
\begin{tcolorbox}[colback=lightbg,colframe=bordergray,left=6pt,right=6pt,top=3pt,bottom=3pt,boxrule=0.6pt,arc=2.5pt]
\small
\noindent\textbf{\color{navy}교수학적 강조점 (Focus Point):}\enspace 선형 ORM의 72개 파라미터와 MLP의 11,628개 파라미터 용량을 대조하여 비선형성 포착 능력을 설명합니다.\par\vspace{0.3em}
\noindent\textbf{\color{successgreen}예상 질문 및 모범 답변 (Q\&A):}\enspace Q: 은닉층을 (32, 16)처럼 더 작게 줄여도 되나요? A: 좋은 질문입니다. 실제로 적절히 설계된 작은 Tanh 신경망이 과적합 없이 더 빠르고 우수한 성능을 내기도 합니다.\par\vspace{0.3em}
\noindent\textbf{\color{accent}다음 슬라이드 연결 멘트 (Transition):}\enspace 그렇다면 이 신경망 대리 모델을 한 번도 본 적 없는 독립 검증 데이터에 평가했을 때 오차는 어느 정도일까요?
\end{tcolorbox}
\clearpage
\hypertarget{slide.20}{}
\phantomsection
\addcontentsline{toc}{section}{슬라이드 20: 대리 모델 훈련 및 독립 홀드아웃(Holdout) 검증 결과}
\noindent\begin{tcolorbox}[colback=navy,colframe=navy,left=6pt,right=6pt,top=3pt,bottom=3pt,boxrule=0pt,arc=2pt]
{\large\bfseries\color{white}슬라이드 20:\enspace 대리 모델 훈련 및 독립 홀드아웃(Holdout) 검증 결과}\hfill{\footnotesize\bfseries\color{white!85!lightbg}Paradigm 1: Supervised Learning \enspace$|$\enspace 슬라이드 20 / 35}
\end{tcolorbox}
\vspace{0.2em}
\begin{tcolorbox}[colback=lightbg,colframe=bordergray,title={\footnotesize\bfseries\color{darkgray}슬라이드 시각 구성 요약},fontupper=\small,left=6pt,right=6pt,top=2.5pt,bottom=2.5pt,boxrule=0.6pt,arc=2.5pt]
데이터셋 요약(1,500개 샷, 80/20 분할)과 홀드아웃 테스트 RMSE: 선형 릿지 0.082 mm, 기준 MLP 0.109 mm, 튜닝된 MLP 0.012 mm, 센서 노이즈 바닥 0.010 mm.
\end{tcolorbox}
\vspace{0.2em}
\begin{tcolorbox}[colback=white,colframe=primaryblue,title={\footnotesize\bfseries\color{white}발표자 구두 대본 (Spoken Script)},fontupper=\small,left=6pt,right=6pt,top=3pt,bottom=3pt,boxrule=0.8pt,arc=2.5pt]
대리 모델이 물리를 제대로 학습했는지 어떻게 알 수 있을까요? 한 번도 본 적 없는 데이터로 평가하면 됩니다. 전류를 [-3.0 A, +3.0 A] 범위에서 균일하게 추출한 1,500개의 오프라인 샷을 만들고, 그중 1,200개로 학습한 뒤 따로 떼어 둔 300개 샷으로 평가합니다. 세 가지 결과를 보겠습니다. 파라미터가 84개인 선형 릿지(Ridge) 모델의 테스트 RMSE는 0.082 mm입니다. 3차 자기 포화를 표현할 수 없기 때문입니다. 파라미터가 11,628개이고 ReLU 활성화 함수와 Adam 옵티마이저를 쓰는 기준 MLP의 테스트 RMSE는 0.109 mm입니다. 선형 모델보다 오히려 나쁩니다! 손실이 더 이상 줄지 않아 Adam이 약 85 에포크 만에 학습을 멈추면서 신경망이 충분히 학습되지 못한 것입니다. 세 번째 모델은 64개 유닛 2개 층에 tanh 활성화 함수를 쓰고 전체 배치 방식의 L-BFGS 솔버로 학습한 더 작은 신경망입니다. 이 모델은 0.012 mm에 도달하여 BPM 노이즈 바닥인 0.010 mm에 근접합니다. 노이즈가 섞인 테스트 데이터에서는 어떤 모델도 이 바닥보다 좋아질 수 없습니다. 교훈은 큰 신경망이 저절로 좋은 대리 모델이 되지는 않는다는 것입니다. 반드시 홀드아웃 오차를 확인해야 하며, 이를 개선하는 것이 바로 TODO 1의 과제입니다.
\end{tcolorbox}
\vspace{0.2em}
\begin{tcolorbox}[colback=lightbg,colframe=bordergray,left=6pt,right=6pt,top=3pt,bottom=3pt,boxrule=0.6pt,arc=2.5pt]
\small
\noindent\textbf{\color{navy}교수학적 강조점 (Focus Point):}\enspace 기본 신경망이 선형 모델보다 나쁘다는 의외의 결과에서 잠시 멈추고 강조합니다. 이어서 0.012 mm가 어떤 모델도 넘을 수 없는 한계인 0.010 mm 노이즈 바닥에 가깝다는 점을 보여줍니다.\par\vspace{0.3em}
\noindent\textbf{\color{successgreen}예상 질문 및 모범 답변 (Q\&A):}\enspace Q: 기준 MLP가 이렇게 부정확한데 베이스라인 1은 어떻게 779점을 얻나요? A: 제어 루프가 닫혀 있기 때문입니다. 매 스텝 새로운 측정값에서 출발하므로 모델 오차가 단계적으로 보정됩니다. 정확한 모델은 주로 안착 정밀도를 높여 줍니다.\par\vspace{0.3em}
\noindent\textbf{\color{accent}다음 슬라이드 연결 멘트 (Transition):}\enspace 순방향 디지털 트윈이 준비되었다면, 이를 폐루프 궤도 제어에 어떻게 활용하는지 살펴보겠습니다.
\end{tcolorbox}
\clearpage
\hypertarget{slide.21}{}
\phantomsection
\addcontentsline{toc}{section}{슬라이드 21: 폐루프 제어 아키텍처 및 대리 모델 역산 흐름}
\noindent\begin{tcolorbox}[colback=navy,colframe=navy,left=6pt,right=6pt,top=3pt,bottom=3pt,boxrule=0pt,arc=2pt]
{\large\bfseries\color{white}슬라이드 21:\enspace 폐루프 제어 아키텍처 및 대리 모델 역산 흐름}\hfill{\footnotesize\bfseries\color{white!85!lightbg}Paradigm 1: Supervised Learning \enspace$|$\enspace 슬라이드 21 / 35}
\end{tcolorbox}
\vspace{0.2em}
\begin{tcolorbox}[colback=lightbg,colframe=bordergray,title={\footnotesize\bfseries\color{darkgray}슬라이드 시각 구성 요약},fontupper=\small,left=6pt,right=6pt,top=2.5pt,bottom=2.5pt,boxrule=0.6pt,arc=2.5pt]
폐루프 피드백 블록도: 가속기 전송선 환경, BPM 측정값 Y, 현재 누적 전류 I, 디지털 트윈 역산 최적화기, 제어 킥 ΔI, 폐루프 피드백 순환 구조.
\end{tcolorbox}
\vspace{0.2em}
\begin{tcolorbox}[colback=white,colframe=primaryblue,title={\footnotesize\bfseries\color{white}발표자 구두 대본 (Spoken Script)},fontupper=\small,left=6pt,right=6pt,top=3pt,bottom=3pt,boxrule=0.8pt,arc=2.5pt]
이 다이어그램은 폐루프 궤도 제어 시스템의 전체 정보 흐름을 보여줍니다. 스텝 t에서 빔라인 BPM 센서로부터 12차원 관측 궤도 Y가 측정됩니다. 이때 제어기는 환경 메타데이터로부터 현재 전자석들에 인가되어 있는 누적 전류 벡터 I\_t(R$^6$)를 함께 읽어옵니다. 이 현재 전류 정보는 대단히 중요합니다! 자기 포화는 누적 전류의 3승에 비례하므로, 현재 자신이 B-H 곡선의 어느 지점에 위치해 있는지를 알지 못하면 올바른 증분 킥(ΔI)을 계산할 수 없습니다. 제어기는 관측 궤도 Y와 현재 전류 I\_t를 모델 역산 최적화기에 입력합니다. 학습된 신경망 디지털 트윈 f\_θ를 활용하여, 다음 스텝 t+1에서 잔여 궤도를 0으로 상쇄시킬 수 있는 최적의 델타 전류 ΔI ∈ [-0.8A, 0.8A]를 계산합니다. 이 킥이 전원 장치에 인가되면 자장이 변경되고, 빔이 편향되어 새로운 BPM 측정값이 관측됩니다. 이 폐루프 루틴은 궤도 RMS 오차가 목표치인 0.030 mm 이하로 떨어질 때까지 반복됩니다.
\end{tcolorbox}
\vspace{0.2em}
\begin{tcolorbox}[colback=lightbg,colframe=bordergray,left=6pt,right=6pt,top=3pt,bottom=3pt,boxrule=0.6pt,arc=2.5pt]
\small
\noindent\textbf{\color{navy}교수학적 강조점 (Focus Point):}\enspace 환경의 info['currents']로부터 현재 누적 전류를 읽어오는 것이 포화 보상의 핵심 열쇠임을 힘주어 설명합니다.\par\vspace{0.3em}
\noindent\textbf{\color{successgreen}예상 질문 및 모범 답변 (Q\&A):}\enspace Q: 현재 전류를 추적하지 않으면 어떻게 되나요? A: 항상 0A 부근의 선형 영역에 있다고 착각하여, 고전류 포화 영역에서 킥을 크게 과소평가하게 됩니다.\par\vspace{0.3em}
\noindent\textbf{\color{accent}다음 슬라이드 연결 멘트 (Transition):}\enspace 이제 수십 년간 가속기 제어의 표준으로 군림해 온 고전 선형 SVD 제어 방식(Baseline 0)과 비교해보겠습니다.
\end{tcolorbox}
\clearpage
\hypertarget{slide.22}{}
\phantomsection
\addcontentsline{toc}{section}{슬라이드 22: 베이스라인 0: 고전적 선형 SVD 의사역행렬 제어기}
\noindent\begin{tcolorbox}[colback=navy,colframe=navy,left=6pt,right=6pt,top=3pt,bottom=3pt,boxrule=0pt,arc=2pt]
{\large\bfseries\color{white}슬라이드 22:\enspace 베이스라인 0: 고전적 선형 SVD 의사역행렬 제어기}\hfill{\footnotesize\bfseries\color{white!85!lightbg}Baseline 0: Classical Linear SVD \enspace$|$\enspace 슬라이드 22 / 35}
\end{tcolorbox}
\vspace{0.2em}
\begin{tcolorbox}[colback=lightbg,colframe=bordergray,title={\footnotesize\bfseries\color{darkgray}슬라이드 시각 구성 요약},fontupper=\small,left=6pt,right=6pt,top=2.5pt,bottom=2.5pt,boxrule=0.6pt,arc=2.5pt]
응답행렬 특이값 분해 R$^T$ = UΣV$^T$, 무어-펜로즈 의사역행렬 (R$^T$)†, 피드백 수식 ΔI = -α·(R$^T$)†·Y, 피드백 게인 α에 따른 수렴 속도 비교.
\end{tcolorbox}
\vspace{0.2em}
\begin{tcolorbox}[colback=white,colframe=primaryblue,title={\footnotesize\bfseries\color{white}발표자 구두 대본 (Spoken Script)},fontupper=\small,left=6pt,right=6pt,top=3pt,bottom=3pt,boxrule=0.8pt,arc=2.5pt]
베이스라인 0은 전통적인 가속기 궤도 보정의 표준인 '특이값 분해(SVD) 기반 무어-펜로즈 의사역행렬 제어기'입니다. 이상적인 선형 모델에서 궤도 변위는 Y = ΔI · R로 표현되며, 전치하면 R$^T$·ΔI = Y가 됩니다. R$^T$은 12×6 행렬이므로 행이 열보다 많은 과결정 시스템입니다. SVD 분해 R$^T$ = UΣV$^T$를 통해 의사역행렬 (R$^T$)† = VΣ$^{-1}$U$^T$를 계산하고, 비례 게인 α를 곱해 제어 킥을 인가합니다: ΔI = -α · (R$^T$)† · Y. 기본 보수적 게인 α = 0.45를 적용하면 30개 시드 전체에서 안정적으로 100\% 성공률을 달성합니다. 하지만 이 제어기는 3차 자기 포화를 전혀 알지 못하므로, 고전류 영역에서 킥을 계속 과소평가하여 수렴까지 평균 9.4스텝이나 소요되며 공식 벤치마크 점수는 635.277점에 머무릅니다. 게인을 α = 1.2로 높이면 약 3.7스텝 만에 수렴하여 822점을 얻지만, 그보다 크면 오버슈트가 발생하여 점수가 다시 떨어집니다.
\end{tcolorbox}
\vspace{0.2em}
\begin{tcolorbox}[colback=lightbg,colframe=bordergray,left=6pt,right=6pt,top=3pt,bottom=3pt,boxrule=0.6pt,arc=2.5pt]
\small
\noindent\textbf{\color{navy}교수학적 강조점 (Focus Point):}\enspace 게인 트레이드오프를 설명합니다: α=0.45는 안전하지만 9.4스텝으로 느리고, α=1.2는 3.7스텝(822점)으로 빠르지만, 그 이상에서는 오버슈트가 발생합니다.\par\vspace{0.3em}
\noindent\textbf{\color{successgreen}예상 질문 및 모범 답변 (Q\&A):}\enspace Q: 왜 R을 전치(transpose)하여 역행렬을 구하나요? A: 수식 정의상 Y가 (1×12) 행 벡터이고 전류 I가 (1×6) 행 벡터이므로(Y = I·R), 전치해야 올바른 차원이 맞습니다.\par\vspace{0.3em}
\noindent\textbf{\color{accent}다음 슬라이드 연결 멘트 (Transition):}\enspace 그렇다면 선형 SVD 대신 우리가 학습시킨 비선형 신경망 대리 모델을 역산에 적용하면 어떤 결과가 나올까요? 베이스라인 1을 살펴보겠습니다.
\end{tcolorbox}
\clearpage
\hypertarget{slide.23}{}
\phantomsection
\addcontentsline{toc}{section}{슬라이드 23: 베이스라인 1 (트랙 1): 딥러닝 대리 모델 역산 최적화}
\noindent\begin{tcolorbox}[colback=navy,colframe=navy,left=6pt,right=6pt,top=3pt,bottom=3pt,boxrule=0pt,arc=2pt]
{\large\bfseries\color{white}슬라이드 23:\enspace 베이스라인 1 (트랙 1): 딥러닝 대리 모델 역산 최적화}\hfill{\footnotesize\bfseries\color{white!85!lightbg}Baseline 1 (Track 1) \enspace$|$\enspace 슬라이드 23 / 35}
\end{tcolorbox}
\vspace{0.2em}
\begin{tcolorbox}[colback=lightbg,colframe=bordergray,title={\footnotesize\bfseries\color{darkgray}슬라이드 시각 구성 요약},fontupper=\small,left=6pt,right=6pt,top=2.5pt,bottom=2.5pt,boxrule=0.6pt,arc=2.5pt]
모델 역산 목적함수 수식: 잔여 궤도 최소화 및 액추에이터 페널티 정규화, SciPy L-BFGS-B 솔버, 평균 4.0스텝 수렴(779.169점).
\end{tcolorbox}
\vspace{0.2em}
\begin{tcolorbox}[colback=white,colframe=primaryblue,title={\footnotesize\bfseries\color{white}발표자 구두 대본 (Spoken Script)},fontupper=\small,left=6pt,right=6pt,top=3pt,bottom=3pt,boxrule=0.8pt,arc=2.5pt]
베이스라인 1은 해커톤 트랙 1의 핵심인 '디지털 트윈 역산 기반 모델 최적화 제어기'입니다. 매 스텝마다 고정된 선형 행렬 대신, 우리가 사전에 학습시킨 MLP 대리 모델을 활용하여 다음과 같은 최적화 문제를 풉니다: min\_{ΔI} || Y - f\_θ(I) + f\_θ(I + ΔI) ||$^2$ + λ ||ΔI||$^2$. 이 수식의 의미를 보십시오: f\_θ(I + ΔI) - f\_θ(I)는 델타 킥 ΔI에 의해 유도될 궤도 변화량을 의미하며, 현재 관측값 Y에 더해주면 다음 스텝의 '예측 잔여 궤도'가 됩니다. 대리 모델 f\_θ가 3차 자기 포화를 정확히 표현할수록, 포화를 극복하기 위해 전류를 얼마나 더 주어야 하는지 최적화 알고리즘이 더 잘 계산해냅니다. SciPy의 L-BFGS-B 솔버로 매 스텝 약 27밀리초 만에 해를 구하고, 킥은 [-0.8A, 0.8A] 스텝 한계로 클리핑합니다. 그 결과 30개 시드에서 평균 4.0스텝 만에 수렴하여 779.169점을 기록하며, 고전 SVD보다 144점 높습니다. 테스트 오차가 0.109 mm인 기준 대리 모델로 얻은 결과라는 점에 주목하십시오. 피드백 루프가 모델 오차의 상당 부분을 보정해 주기 때문입니다.
\end{tcolorbox}
\vspace{0.2em}
\begin{tcolorbox}[colback=lightbg,colframe=bordergray,left=6pt,right=6pt,top=3pt,bottom=3pt,boxrule=0.6pt,arc=2.5pt]
\small
\noindent\textbf{\color{navy}교수학적 강조점 (Focus Point):}\enspace 635.277점(SVD)에서 779.169점(베이스라인 1)으로의 점수 상승이 모델 기반의 더 큰 스텝에서 나온다는 점을 짚어 줍니다. 대리 모델이 더 정확해지면 정밀도 점수가 추가로 올라갑니다.\par\vspace{0.3em}
\noindent\textbf{\color{successgreen}예상 질문 및 모범 답변 (Q\&A):}\enspace Q: 트랙 1에서 779점보다 더 높은 점수를 낼 수 있나요? A: 네. TODO 1에서 대리 모델 정밀도를 높이면(예: tanh와 lbfgs, 0.012 mm) 같은 제어기로 약 900점, 주최 측 솔루션 제어기로는 924점을 얻습니다.\par\vspace{0.3em}
\noindent\textbf{\color{accent}다음 슬라이드 연결 멘트 (Transition):}\enspace 이어서 최적화 역산 대신 실시간 피드백 정책을 학습하는 두 번째 패러다임, 강화학습(Reinforcement Learning)으로 넘어가겠습니다.
\end{tcolorbox}
\clearpage
\hypertarget{slide.24}{}
\phantomsection
\addcontentsline{toc}{section}{슬라이드 24: 강화학습: 마르코프 결정 과정(MDP) 정식화 및 하이퍼파라미터}
\noindent\begin{tcolorbox}[colback=navy,colframe=navy,left=6pt,right=6pt,top=3pt,bottom=3pt,boxrule=0pt,arc=2pt]
{\large\bfseries\color{white}슬라이드 24:\enspace 강화학습: 마르코프 결정 과정(MDP) 정식화 및 하이퍼파라미터}\hfill{\footnotesize\bfseries\color{white!85!lightbg}Paradigm 2: Reinforcement Learning \enspace$|$\enspace 슬라이드 24 / 35}
\end{tcolorbox}
\vspace{0.2em}
\begin{tcolorbox}[colback=lightbg,colframe=bordergray,title={\footnotesize\bfseries\color{darkgray}슬라이드 시각 구성 요약},fontupper=\small,left=6pt,right=6pt,top=2.5pt,bottom=2.5pt,boxrule=0.6pt,arc=2.5pt]
좌측: 폐루프 MDP 프레임워크(상태, 행동, 전이, 보상 함수, 할인율), 누적 보상 J(θ). 우측: 참가자 튜닝 변수(TODO 2)가 강조된 RL 하이퍼파라미터 분류표.
\end{tcolorbox}
\vspace{0.2em}
\begin{tcolorbox}[colback=white,colframe=primaryblue,title={\footnotesize\bfseries\color{white}발표자 구두 대본 (Spoken Script)},fontupper=\small,left=6pt,right=6pt,top=3pt,bottom=3pt,boxrule=0.8pt,arc=2.5pt]
강화학습(RL)은 궤도 보정을 마르코프 결정 과정(MDP) 기반의 순차적 의사결정 문제로 다룹니다. 좌측을 보시면, 상태 공간 S는 12개 BPM의 변위 측정값 Y입니다. 행동 공간 A는 6개 스티어러에 인가할 전류 증분 ΔI ∈ [-0.8A, 0.8A]$^6$입니다. 전이 확률 P는 가속기 빔 물리, 자기 포화, 센서 노이즈가 결합된 환경 동역학입니다. 보상 함수 R은 가속기 운전 목표를 정밀하게 반영합니다: r\_t = - w\_rms·RMS\_{t+1} - w\_act·||a\_t||$^2$ - c\_step + b\_term. 잔여 궤도 오차 페널티, 코일 발열 억제를 위한 액션 페널티, 빠른 수렴을 유도하는 스텝 페널티(c\_step), 그리고 목표 달성 보너스(b\_term)로 구성됩니다. 우측 표에는 노트북 TODO 2에서 수정 가능한 변수들이 정리되어 있습니다: 할인율 γ(0.95~0.99), 정책 학습률, 하모닉 댐핑 계수 κ, 보상 가중치(w\_rms, w\_act), 스텝 페널티(c\_step) 등이 포함됩니다.
\end{tcolorbox}
\vspace{0.2em}
\begin{tcolorbox}[colback=lightbg,colframe=bordergray,left=6pt,right=6pt,top=3pt,bottom=3pt,boxrule=0.6pt,arc=2.5pt]
\small
\noindent\textbf{\color{navy}교수학적 강조점 (Focus Point):}\enspace 스텝 페널티(c\_step)의 역할을 강조합니다: 매 스텝마다 감점을 부여함으로써 에이전트가 최소 스텝으로 수렴하도록 유도합니다.\par\vspace{0.3em}
\noindent\textbf{\color{successgreen}예상 질문 및 모범 답변 (Q\&A):}\enspace Q: 유한한 에피소드에서도 할인율 γ < 1이 필요한가요? A: 네, 미래 보상을 적절히 할인함으로써 당장 궤도를 줄이는 즉각적 행동에 높은 가중치를 부여하기 위함입니다.\par\vspace{0.3em}
\noindent\textbf{\color{accent}다음 슬라이드 연결 멘트 (Transition):}\enspace 그렇다면 위험한 무작위 탐색 없이, 안정적이면서도 빠르게 동작하는 강화학습 제어기를 어떻게 구성할 수 있을까요? 트랙 2를 보시겠습니다.
\end{tcolorbox}
\clearpage
\hypertarget{slide.25}{}
\phantomsection
\addcontentsline{toc}{section}{슬라이드 25: 트랙 2: 물리 사전 지식 결합 폐루프 강화학습 정책}
\noindent\begin{tcolorbox}[colback=navy,colframe=navy,left=6pt,right=6pt,top=3pt,bottom=3pt,boxrule=0pt,arc=2pt]
{\large\bfseries\color{white}슬라이드 25:\enspace 트랙 2: 물리 사전 지식 결합 폐루프 강화학습 정책}\hfill{\footnotesize\bfseries\color{white!85!lightbg}Track 2 Challenge \enspace$|$\enspace 슬라이드 25 / 35}
\end{tcolorbox}
\vspace{0.2em}
\begin{tcolorbox}[colback=lightbg,colframe=bordergray,title={\footnotesize\bfseries\color{darkgray}슬라이드 시각 구성 요약},fontupper=\small,left=6pt,right=6pt,top=2.5pt,bottom=2.5pt,boxrule=0.6pt,arc=2.5pt]
물리 사전 지식 기반 잔차 정책 구조도: ΔI = γ(t)·[ -α·(R$^T$)†·Y + W\_learned·Y + b ], 해석적 선형 사전 지식과 학습된 비선형 보정 잔차의 결합.
\end{tcolorbox}
\vspace{0.2em}
\begin{tcolorbox}[colback=white,colframe=primaryblue,title={\footnotesize\bfseries\color{white}발표자 구두 대본 (Spoken Script)},fontupper=\small,left=6pt,right=6pt,top=3pt,bottom=3pt,boxrule=0.8pt,arc=2.5pt]
가속기 환경에서 아무런 사전 지식 없이 바닥부터(tabula rasa) 강화학습을 수행하면, 초기에 무작위 킥을 인가하여 빔이 챔버에 충돌하게 됩니다. 트랙 2에서는 이를 해결하기 위해 '물리 사전 지식 결합 잔차 정책(Physics-Informed Residual Policy)'을 도입했습니다. 제어 방정식을 보십시오: ΔI(t) = γ(t) · [ -α·(R$^T$)†·Y + W\_learned·Y + b ]. 앞부분의 해석적 물리 사전 지식은 우리가 잘 알고 있는 SVD 의사역행렬입니다. 궤도 보정의 대부분을 즉각적이고 안정적으로 처리합니다. 뒷부분의 학습된 잔차(W\_learned, b)는 선형 모델이 설명하지 못하는 '3차 자기 포화 왜곡'만을 전담하여 학습합니다! 기본 제공되는 기준 제어기에서는 W\_learned와 b가 0으로 초기화되어 있어 단순 감쇠 SVD로 동작하며 590.559점을 기록합니다. 낮은 게인 0.45에서는 감쇠가 수렴을 늦추기만 하므로 베이스라인 0보다 낮은 점수입니다. 트랙 2에서 참가자 여러분의 과제는 이 잔차 파라미터를 학습시키거나 보정하여 비선형 포화를 완벽히 극복하는 고성능 정책을 완성하는 것입니다.
\end{tcolorbox}
\vspace{0.2em}
\begin{tcolorbox}[colback=lightbg,colframe=bordergray,left=6pt,right=6pt,top=3pt,bottom=3pt,boxrule=0.6pt,arc=2.5pt]
\small
\noindent\textbf{\color{navy}교수학적 강조점 (Focus Point):}\enspace 우수한 물리 사전 지식 주변의 '잔차(residual)'만을 학습하는 기법이 바닥부터 학습하는 것보다 훨씬 쉽고 안전함을 설명합니다.\par\vspace{0.3em}
\noindent\textbf{\color{successgreen}예상 질문 및 모범 답변 (Q\&A):}\enspace Q: 선형 SVD 대신 비선형 역산 모델을 사전 지식으로 교체해도 되나요? A: 완벽한 전략입니다! 주최 측 솔루션에서는 포화 보상 사전 지식을 적용하여 939.981점을 달성했습니다.\par\vspace{0.3em}
\noindent\textbf{\color{accent}다음 슬라이드 연결 멘트 (Transition):}\enspace 정책 수식 맨 앞에 곱해진 감쇠 계수 γ(t)는 과연 어떤 역할을 수행할까요? 슬라이드 26에서 그 비밀을 밝혀보겠습니다.
\end{tcolorbox}
\clearpage
\hypertarget{slide.26}{}
\phantomsection
\addcontentsline{toc}{section}{슬라이드 26: 동적 조화 스텝 감쇠(Dynamic Harmonic Step Damping)}
\noindent\begin{tcolorbox}[colback=navy,colframe=navy,left=6pt,right=6pt,top=3pt,bottom=3pt,boxrule=0pt,arc=2pt]
{\large\bfseries\color{white}슬라이드 26:\enspace 동적 조화 스텝 감쇠(Dynamic Harmonic Step Damping)}\hfill{\footnotesize\bfseries\color{white!85!lightbg}Track 2 Challenge \enspace$|$\enspace 슬라이드 26 / 35}
\end{tcolorbox}
\vspace{0.2em}
\begin{tcolorbox}[colback=lightbg,colframe=bordergray,title={\footnotesize\bfseries\color{darkgray}슬라이드 시각 구성 요약},fontupper=\small,left=6pt,right=6pt,top=2.5pt,bottom=2.5pt,boxrule=0.6pt,arc=2.5pt]
감쇠 스케줄 수식 γ(t) = 1.0 / (1.0 + κ·t), 그리고 게인 0.45, 1.0, 1.2에서 감쇠 유무에 따라 측정한 벤치마크 점수.
\end{tcolorbox}
\vspace{0.2em}
\begin{tcolorbox}[colback=white,colframe=primaryblue,title={\footnotesize\bfseries\color{white}발표자 구두 대본 (Spoken Script)},fontupper=\small,left=6pt,right=6pt,top=3pt,bottom=3pt,boxrule=0.8pt,arc=2.5pt]
제어기가 0.030 mm 목표에 가까워지면 측정된 궤도 오차가 BPM의 10 ㎛ 노이즈와 비슷한 크기가 됩니다. 이때 게인이 높은 제어기는 노이즈에도 반응하여, 안착하지 못하고 목표 경계를 넘나들 수 있습니다. 이를 완화하는 기법이 '동적 조화 스텝 감쇠'입니다: γ(t) = 1.0 / (1.0 + κ·t), 기본 감쇠 계수 κ = 0.03. 스텝 0에서는 γ가 정확히 1.0이므로 큰 초기 보정에 제어 권한을 모두 사용합니다. 이후 완만하게 감소하여 10스텝에서 약 0.77, 30스텝에서 약 0.53이 됩니다. 그러나 감쇠는 공짜 이득이 아니라 트레이드오프이며, 슬라이드의 수치가 이를 보여줍니다. 게인 0.45의 단순 의사역행렬 제어기에서는 감쇠가 점수를 635점에서 591점으로 낮춥니다. 애초에 노이즈 문제가 없었으므로 스텝만 늘어난 것입니다. 게인 1.0에서는 765점에서 783점으로 올라갑니다. 더 안정적인 안착이 추가 스텝보다 가치가 크기 때문입니다. 게인 1.2에서는 822점에서 809점으로 다시 낮아집니다. 그리고 강한 감쇠(κ = 0.1)는 모든 경우에 손해입니다. 우리 벤치마크에서는 처음으로 30 ㎛ 아래 측정값이 나오는 순간 에피소드가 끝나므로, 몇 시간 동안 궤도를 유지해야 하는 실제 피드백 시스템보다 감쇠의 중요성이 작습니다. κ는 규칙이 아니라 시험해 볼 파라미터로 다루십시오.
\end{tcolorbox}
\vspace{0.2em}
\begin{tcolorbox}[colback=lightbg,colframe=bordergray,left=6pt,right=6pt,top=3pt,bottom=3pt,boxrule=0.6pt,arc=2.5pt]
\small
\noindent\textbf{\color{navy}교수학적 강조점 (Focus Point):}\enspace 감쇠가 트레이드오프임을 강조합니다. 게인 0.45와 1.2에서는 점수를 잃었고 게인 1.0에서는 점수를 얻었습니다. 가정하지 말고 직접 측정해 보도록 권합니다.\par\vspace{0.3em}
\noindent\textbf{\color{successgreen}예상 질문 및 모범 답변 (Q\&A):}\enspace Q: κ 값이 0.2처럼 너무 크면 어떻게 되나요? A: 제어 권한이 너무 빨리 줄어들어 훨씬 많은 스텝이 필요해집니다. κ = 0.1만 되어도 시험한 모든 게인에서 점수가 낮아졌습니다.\par\vspace{0.3em}
\noindent\textbf{\color{accent}다음 슬라이드 연결 멘트 (Transition):}\enspace 이제 세 번째 패러다임이자 샘플 효율성의 정점인 가우시안 프로세스 기반 베이지안 최적화(트랙 3)를 살펴보겠습니다.
\end{tcolorbox}
\clearpage
\hypertarget{slide.27}{}
\phantomsection
\addcontentsline{toc}{section}{슬라이드 27: 베이지안 최적화: 확률론적 가우시안 프로세스(GP) 및 하이퍼파라미터}
\noindent\begin{tcolorbox}[colback=navy,colframe=navy,left=6pt,right=6pt,top=3pt,bottom=3pt,boxrule=0pt,arc=2pt]
{\large\bfseries\color{white}슬라이드 27:\enspace 베이지안 최적화: 확률론적 가우시안 프로세스(GP) 및 하이퍼파라미터}\hfill{\footnotesize\bfseries\color{white!85!lightbg}Paradigm 3: Bayesian Optimization \enspace$|$\enspace 슬라이드 27 / 35}
\end{tcolorbox}
\vspace{0.2em}
\begin{tcolorbox}[colback=lightbg,colframe=bordergray,title={\footnotesize\bfseries\color{darkgray}슬라이드 시각 구성 요약},fontupper=\small,left=6pt,right=6pt,top=2.5pt,bottom=2.5pt,boxrule=0.6pt,arc=2.5pt]
좌측: 베이지안 최적화 4단계 흐름(비모수 사전분포, 사후 조건부 분포 도출, 획득함수 최적화, 빔라인 측정 갱신). 우측: 참가자 튜닝 변수(TODO 3)가 강조된 BO 하이퍼파라미터 분류표.
\end{tcolorbox}
\vspace{0.2em}
\begin{tcolorbox}[colback=white,colframe=primaryblue,title={\footnotesize\bfseries\color{white}발표자 구두 대본 (Spoken Script)},fontupper=\small,left=6pt,right=6pt,top=3pt,bottom=3pt,boxrule=0.8pt,arc=2.5pt]
베이지안 최적화(BO)는 샘플 효율성이 극대화된 블랙박스 최적화 기법입니다. 수천 개의 데이터가 필요한 신경망과 달리, 가우시안 프로세스(GP)는 수십 개의 샷만으로도 물리 함수를 정확히 모델링하며 각 지점의 '불확실성(uncertainty)'을 정량적으로 제공합니다. 좌측의 4단계 워크플로우를 보십시오: 1단계에서는 비모수적 GP 사전분포 f(x) ~ GP(m, k + σ\_n$^2$)를 설정합니다. 2단계에서는 측정된 데이터셋 D\_t에 대한 조건부 확률을 해석적으로 계산하여 예측 평균 μ(x)과 불확실성 σ(x)을 도출합니다. 3단계에서는 하한 신뢰 구간(LCB) 획득 함수를 최적화하여 탐색(exploration)과 활용(exploitation)의 균형점을 찾습니다. 4단계에서는 도출된 최적 킥을 빔라인에 인가하고 데이터를 갱신합니다. 우측 표에는 노트북 TODO 3에서 튜닝할 수 있는 변수들이 정리되어 있습니다: 마테른(Matérn-5/2) 커널, 자석별 상관 길이(lengthscale) ℓ, 신호 분산 σ\_f$^2$, 고정 노이즈 바닥 σ\_n = 0.010 mm, 탐색 파라미터 κ, 정규화 계수 λ 등이 포함됩니다.
\end{tcolorbox}
\vspace{0.2em}
\begin{tcolorbox}[colback=lightbg,colframe=bordergray,left=6pt,right=6pt,top=3pt,bottom=3pt,boxrule=0.6pt,arc=2.5pt]
\small
\noindent\textbf{\color{navy}교수학적 강조점 (Focus Point):}\enspace 노이즈 분산 σ\_n = 0.010 mm는 물리적 센서 특성이므로 절대 수정해서는 안 되는 고정 파라미터임을 명시합니다.\par\vspace{0.3em}
\noindent\textbf{\color{successgreen}예상 질문 및 모범 답변 (Q\&A):}\enspace Q: 왜 RBF 커널 대신 Matérn-5/2를 쓰나요? A: RBF는 무한 번 미분 가능한 비현실적인 매끄러움을 가정하지만, Matérn-5/2는 2회 미분 가능하여 실제 전자석 포화 특성에 완벽히 부합합니다.\par\vspace{0.3em}
\noindent\textbf{\color{accent}다음 슬라이드 연결 멘트 (Transition):}\enspace 이 가우시안 프로세스 모델이 해커톤 트랙 3 폐루프 제어기 내부에서 어떻게 구동되는지 살펴보겠습니다.
\end{tcolorbox}
\clearpage
\hypertarget{slide.28}{}
\phantomsection
\addcontentsline{toc}{section}{슬라이드 28: 트랙 3: 가우시안 프로세스 기반 폐루프 베이지안 최적화}
\noindent\begin{tcolorbox}[colback=navy,colframe=navy,left=6pt,right=6pt,top=3pt,bottom=3pt,boxrule=0pt,arc=2pt]
{\large\bfseries\color{white}슬라이드 28:\enspace 트랙 3: 가우시안 프로세스 기반 폐루프 베이지안 최적화}\hfill{\footnotesize\bfseries\color{white!85!lightbg}Track 3 Challenge \enspace$|$\enspace 슬라이드 28 / 35}
\end{tcolorbox}
\vspace{0.2em}
\begin{tcolorbox}[colback=lightbg,colframe=bordergray,title={\footnotesize\bfseries\color{darkgray}슬라이드 시각 구성 요약},fontupper=\small,left=6pt,right=6pt,top=2.5pt,bottom=2.5pt,boxrule=0.6pt,arc=2.5pt]
트랙 3 제어기 아키텍처, 300회 오프라인 샷 기반 GP 캘리브레이션, 폐루프 스텝별 LCB 최적화, 콜백을 활용한 실시간 온라인 GP 갱신 확장성.
\end{tcolorbox}
\vspace{0.2em}
\begin{tcolorbox}[colback=white,colframe=primaryblue,title={\footnotesize\bfseries\color{white}발표자 구두 대본 (Spoken Script)},fontupper=\small,left=6pt,right=6pt,top=3pt,bottom=3pt,boxrule=0.8pt,arc=2.5pt]
트랙 3 제어기의 구동 원리입니다. 기준 코드에서는 [-3.0A, +3.0A] 범위에서 수집된 300개의 캘리브레이션 샷으로 Matérn-5/2 커널 GP를 사전 학습시킵니다. 매 스텝마다 제어기는 최신 BPM 관측값 Y와 현재 누적 전류 I\_t를 전달받습니다. 대리 모델을 통해 후보 킥 ΔI에 대한 예측 평균 궤도와 불확실성을 계산하고, LCB 획득 함수를 최소화하는 최적의 델타 전류를 도출합니다. 여기서 매우 중요한 실무적 포인트를 짚어드립니다: 기준 코드의 트랙 3 제어기는 매 스텝 최신 측정값으로부터 최적화를 다시 풀기 때문에 '폐루프(closed-loop)' 제어기이지만, GP 모델 자체는 사전에 고정된 오프라인 모델입니다. 참가자 여러분께서 노트북의 CALLBACKS 딕셔너리에 step\_callback을 등록하여 매 스텝 획득한 새로운 측정값을 GP에 실시간 추가하고 재학습(refit)시킨다면, 진정한 '온라인 실시간 베이지안 최적화'로 업그레이드할 수 있습니다!
\end{tcolorbox}
\vspace{0.2em}
\begin{tcolorbox}[colback=lightbg,colframe=bordergray,left=6pt,right=6pt,top=3pt,bottom=3pt,boxrule=0.6pt,arc=2.5pt]
\small
\noindent\textbf{\color{navy}교수학적 강조점 (Focus Point):}\enspace 폐루프 제어와 온라인 모델 갱신의 차이를 명확히 구분해주고, 콜백을 활용한 온라인 GP 갱신이 고득점의 비결임을 귀띔합니다.\par\vspace{0.3em}
\noindent\textbf{\color{successgreen}예상 질문 및 모범 답변 (Q\&A):}\enspace Q: 실시간으로 GP를 재학습하면 속도가 너무 느려지지 않나요? A: 맞습니다. GP 연산량은 데이터 수의 3승(O(N$^3$))에 비례하므로, 최근 50개 정도의 슬라이딩 윈도우 버퍼만 유지하는 것이 좋습니다.\par\vspace{0.3em}
\noindent\textbf{\color{accent}다음 슬라이드 연결 멘트 (Transition):}\enspace 그렇다면 트랙 3의 핵심 의사결정 수식인 하한 신뢰 구간(LCB) 획득 함수를 자세히 살펴보겠습니다.
\end{tcolorbox}
\clearpage
\hypertarget{slide.29}{}
\phantomsection
\addcontentsline{toc}{section}{슬라이드 29: 하한 신뢰 구간(Lower Confidence Bound, LCB) 획득 함수}
\noindent\begin{tcolorbox}[colback=navy,colframe=navy,left=6pt,right=6pt,top=3pt,bottom=3pt,boxrule=0pt,arc=2pt]
{\large\bfseries\color{white}슬라이드 29:\enspace 하한 신뢰 구간(Lower Confidence Bound, LCB) 획득 함수}\hfill{\footnotesize\bfseries\color{white!85!lightbg}Track 3 Challenge \enspace$|$\enspace 슬라이드 29 / 35}
\end{tcolorbox}
\vspace{0.2em}
\begin{tcolorbox}[colback=lightbg,colframe=bordergray,title={\footnotesize\bfseries\color{darkgray}슬라이드 시각 구성 요약},fontupper=\small,left=6pt,right=6pt,top=2.5pt,bottom=2.5pt,boxrule=0.6pt,arc=2.5pt]
LCB 획득 함수 수식: 예측 RMS 잔여 오차(활용) - κ·불확실성(탐색) + λ·액션 페널티, 탐색 계수 κ = 0일 때의 3.6스텝 초고속 수렴 성과.
\end{tcolorbox}
\vspace{0.2em}
\begin{tcolorbox}[colback=white,colframe=primaryblue,title={\footnotesize\bfseries\color{white}발표자 구두 대본 (Spoken Script)},fontupper=\small,left=6pt,right=6pt,top=3pt,bottom=3pt,boxrule=0.8pt,arc=2.5pt]
트랙 3의 핵심 의사결정 수식인 하한 신뢰 구간(LCB) 획득 함수입니다: α\_LCB(ΔI) = RMS(예측 다음 궤도) - κ · σ̄(불확실성) + λ · ||ΔI||$^2$. 첫 번째 항은 '활용(exploitation)'으로, 예측 잔여 궤도 오차를 0으로 줄이는 역할을 합니다. 두 번째 항은 '탐색(exploration)'으로, 불확실성이 높은 미지의 자장 영역을 탐색하도록 장려합니다. 탐색 가중치 κ가 이 둘의 균형을 결정합니다. 주목할 점은, 우리 해커톤 환경에서는 사전 300회 샷을 통해 GP 모델이 이미 매우 정밀하게 보정되어 있다는 사실입니다. 따라서 κ = 0으로 설정하여 사후 평균만을 순수하게 활용(pure exploitation)했을 때, 전체 30개 시드에서 평균 단 3.6스텝 만에 목표치에 안착하는 놀라운 속도를 기록했습니다! 공식 점수 878.384점을 획득하며 현재 벤치마크 랭킹 1위를 차지하고 있습니다.
\end{tcolorbox}
\vspace{0.2em}
\begin{tcolorbox}[colback=lightbg,colframe=bordergray,left=6pt,right=6pt,top=3pt,bottom=3pt,boxrule=0.6pt,arc=2.5pt]
\small
\noindent\textbf{\color{navy}교수학적 강조점 (Focus Point):}\enspace 0.8A 스텝 한계가 정하는 하한(평균 약 3.7스텝)에 도달한 3.6스텝의 수렴 속도를 강조합니다. 기준 제어기 중 1위입니다.\par\vspace{0.3em}
\noindent\textbf{\color{successgreen}예상 질문 및 모범 답변 (Q\&A):}\enspace Q: κ > 0 탐색이 필요한 경우는 언제인가요? A: 사전에 캘리브레이션되지 않은 낯선 자장 영역으로 진입하거나 미지의 드리프트가 발생했을 때 필수적입니다.\par\vspace{0.3em}
\noindent\textbf{\color{accent}다음 슬라이드 연결 멘트 (Transition):}\enspace 이제 오늘 오후 여러분이 직접 코딩하게 될 Gymnasium 표준 시뮬레이션 환경 규격을 살펴보겠습니다.
\end{tcolorbox}
\clearpage
\hypertarget{slide.30}{}
\phantomsection
\addcontentsline{toc}{section}{슬라이드 30: 체육관(Gymnasium) 환경 API 표준 규격}
\noindent\begin{tcolorbox}[colback=navy,colframe=navy,left=6pt,right=6pt,top=3pt,bottom=3pt,boxrule=0pt,arc=2pt]
{\large\bfseries\color{white}슬라이드 30:\enspace 체육관(Gymnasium) 환경 API 표준 규격}\hfill{\footnotesize\bfseries\color{white!85!lightbg}Gymnasium Environment Contract \enspace$|$\enspace 슬라이드 30 / 35}
\end{tcolorbox}
\vspace{0.2em}
\begin{tcolorbox}[colback=lightbg,colframe=bordergray,title={\footnotesize\bfseries\color{darkgray}슬라이드 시각 구성 요약},fontupper=\small,left=6pt,right=6pt,top=2.5pt,bottom=2.5pt,boxrule=0.6pt,arc=2.5pt]
표준 Gymnasium API 인터페이스: reset(seed) 및 step(action) 입출력 명세, BenchmarkController.act(obs, info) 표준 클래스 규약.
\end{tcolorbox}
\vspace{0.2em}
\begin{tcolorbox}[colback=white,colframe=primaryblue,title={\footnotesize\bfseries\color{white}발표자 구두 대본 (Spoken Script)},fontupper=\small,left=6pt,right=6pt,top=3pt,bottom=3pt,boxrule=0.8pt,arc=2.5pt]
벤치마크의 학술적 엄밀성과 재현성을 위해, 시뮬레이션 환경 AcceleratorOrbitGymEnv는 글로벌 표준인 Farama Gymnasium API를 엄격히 준수합니다. 환경은 두 가지 핵심 메서드를 제공합니다. 첫째, reset(seed): 모든 스티어러 전류를 0으로 리셋하고 해당 시드의 초기 왜곡 Y$_0$를 인가한 뒤, 초기 12차원 BPM 관측값(obs)과 런타임 메타데이터가 담긴 info 딕셔너리를 반환합니다. 둘째, step(action): 6차원 델타 전류 액션을 받아 [-0.8A, 0.8A]로 클리핑한 후 환경을 전진시킵니다. 반환값은 5가지입니다: 새로운 관측값 obs, 보상 r, 성공 종료 여부 terminated(RMS < 0.030 mm 도달 시 True), 최대 스텝 초과 중단 여부 truncated(50스텝 도달 시 True), 그리고 최신 info 딕셔너리입니다. 여러분이 작성할 제어기는 6차원 액션을 반환하는 표준 act(obs, info) 메서드만 구현하면 되며, 특정 클래스를 상속할 필요는 없습니다.
\end{tcolorbox}
\vspace{0.2em}
\begin{tcolorbox}[colback=lightbg,colframe=bordergray,left=6pt,right=6pt,top=3pt,bottom=3pt,boxrule=0.6pt,arc=2.5pt]
\small
\noindent\textbf{\color{navy}교수학적 강조점 (Focus Point):}\enspace terminated=True는 0.03 mm 도달 성공을, truncated=True는 50스텝 타임아웃 실패를 의미함을 명확히 구분 짓습니다.\par\vspace{0.3em}
\noindent\textbf{\color{successgreen}예상 질문 및 모범 답변 (Q\&A):}\enspace Q: act() 메서드 안에서 info['currents']를 사용해도 되나요? A: 네, 적극 권장합니다! 실제 가속기 전원 공급 장치에서도 현재 전류는 실시간 텔레메트리로 제공됩니다.\par\vspace{0.3em}
\noindent\textbf{\color{accent}다음 슬라이드 연결 멘트 (Transition):}\enspace 그렇다면 30개 시드에 대해 제어기 성능은 어떤 공식으로 채점될까요? 슬라이드 31의 공식 채점 수식을 확인해보겠습니다.
\end{tcolorbox}
\clearpage
\hypertarget{slide.31}{}
\phantomsection
\addcontentsline{toc}{section}{슬라이드 31: 공식 30개 시드 평가 프로토콜 및 채점 공식}
\noindent\begin{tcolorbox}[colback=navy,colframe=navy,left=6pt,right=6pt,top=3pt,bottom=3pt,boxrule=0pt,arc=2pt]
{\large\bfseries\color{white}슬라이드 31:\enspace 공식 30개 시드 평가 프로토콜 및 채점 공식}\hfill{\footnotesize\bfseries\color{white!85!lightbg}Evaluation Protocol \enspace$|$\enspace 슬라이드 31 / 35}
\end{tcolorbox}
\vspace{0.2em}
\begin{tcolorbox}[colback=lightbg,colframe=bordergray,title={\footnotesize\bfseries\color{darkgray}슬라이드 시각 구성 요약},fontupper=\small,left=6pt,right=6pt,top=2.5pt,bottom=2.5pt,boxrule=0.6pt,arc=2.5pt]
3대 핵심 채점 공식: 성공 점수(400점) + 속도 점수(최대 300점, 1~20스텝 선형 감쇠) + 정밀도 점수(최대 300점, 30 ㎛에서 10 ㎛ 노이즈 바닥까지 로그 스케일).
\end{tcolorbox}
\vspace{0.2em}
\begin{tcolorbox}[colback=white,colframe=primaryblue,title={\footnotesize\bfseries\color{white}발표자 구두 대본 (Spoken Script)},fontupper=\small,left=6pt,right=6pt,top=3pt,bottom=3pt,boxrule=0.8pt,arc=2.5pt]
공정하고 엄밀한 벤치마크 평가를 위해, 모든 참가자의 제어기는 1번부터 30번까지 30개의 고정 표준 시드에서 평가됩니다. 각 시드 k에서 획득하는 점수 s\_k는 세 가지 기둥에 기반하여 0점에서 1,000점 만점으로 계산됩니다. 만약 50스텝 내에 0.030 mm 목표를 달성하지 못하면 해당 시드는 가차 없이 '0점' 처리됩니다! 시드 하나가 실패하면 그 시드에서 최대 1,000점, 최종 평균으로는 약 33점을 잃으므로 강건성이 최우선입니다. 성공한 시드에 대해서는 세 가지 점수가 합산됩니다: 첫째, '성공 점수(400점)': 0.030 mm 이하 도달 즉시 400점 전액이 부여됩니다. 둘째, '수렴 속도 점수(최대 300점)': 300 × (20 - n\_k) / 19 수식에 따라, 1스텝 만에 수렴하면 300점 만점, 20스텝 이상 소요되면 0점 처리됩니다. 셋째, '안착 정밀도 점수(최대 300점)': 0.030 mm 목표선에서 센서 노이즈 바닥인 0.010 mm까지 로그 스케일로 부여됩니다: 300 × ln(0.030 / RMS) / ln(0.030 / 0.010). 잔여 궤도를 노이즈 바닥에 가깝게 안착시킬수록 정밀도 점수가 만점에 수렴합니다. 30개 시드의 점수를 산술 평균한 값이 여러분의 최종 챌린지 점수가 됩니다.
\end{tcolorbox}
\vspace{0.2em}
\begin{tcolorbox}[colback=lightbg,colframe=bordergray,left=6pt,right=6pt,top=3pt,bottom=3pt,boxrule=0.6pt,arc=2.5pt]
\small
\noindent\textbf{\color{navy}교수학적 강조점 (Focus Point):}\enspace 성공 400점, 속도 300점, 정밀도 300점의 균형을 강조하며, 단 1개 시드의 실패(0점)가 전체 순위에 치명적임을 경고합니다.\par\vspace{0.3em}
\noindent\textbf{\color{successgreen}예상 질문 및 모범 답변 (Q\&A):}\enspace Q: 전류 크기(current norm)는 왜 점수에 반영되지 않나요? A: 보정에 필요한 누적 전류는 시드의 초기 왜곡 크기에 의해 대부분 결정되며, 제어기간 편차가 2\% 미만으로 미미하기 때문입니다.\par\vspace{0.3em}
\noindent\textbf{\color{accent}다음 슬라이드 연결 멘트 (Transition):}\enspace 이 공식 채점 기준에 따라 기존 베이스라인과 챌린지 트랙 제어기들의 성능을 직접 비교해보겠습니다.
\end{tcolorbox}
\clearpage
\hypertarget{slide.32}{}
\phantomsection
\addcontentsline{toc}{section}{슬라이드 32: 공식 벤치마크 리더보드 및 수렴 궤적 비교}
\noindent\begin{tcolorbox}[colback=navy,colframe=navy,left=6pt,right=6pt,top=3pt,bottom=3pt,boxrule=0pt,arc=2pt]
{\large\bfseries\color{white}슬라이드 32:\enspace 공식 벤치마크 리더보드 및 수렴 궤적 비교}\hfill{\footnotesize\bfseries\color{white!85!lightbg}Leaderboard \& Trajectories \enspace$|$\enspace 슬라이드 32 / 35}
\end{tcolorbox}
\vspace{0.2em}
\begin{tcolorbox}[colback=lightbg,colframe=bordergray,title={\footnotesize\bfseries\color{darkgray}슬라이드 시각 구성 요약},fontupper=\small,left=6pt,right=6pt,top=2.5pt,bottom=2.5pt,boxrule=0.6pt,arc=2.5pt]
공식 리더보드 결과표: 1위 베이지안 최적화(878.384점, 3.6스텝), 2위 지도학습 MLP(779.169점, 4.0스텝), 3위 고전 SVD(635.277점, 9.4스텝), 4위 기준 RL(590.559점, 10.8스텝).
\end{tcolorbox}
\vspace{0.2em}
\begin{tcolorbox}[colback=white,colframe=primaryblue,title={\footnotesize\bfseries\color{white}발표자 구두 대본 (Spoken Script)},fontupper=\small,left=6pt,right=6pt,top=3pt,bottom=3pt,boxrule=0.8pt,arc=2.5pt]
30개 표준 평가 시드 전체에 대한 공식 기준 리더보드입니다. 네 가지 제어기 모두 100\% 성공률을 달성했으므로, 순위는 전적으로 수렴 속도와 안착 정밀도에 의해 갈렸습니다. 영예의 1위는 '트랙 3 베이지안 최적화(BO)'입니다: 보정된 GP와 κ = 0 설정을 통해 평균 3.6스텝 만에 수렴하고, 잔여 RMS를 0.0136 mm까지 극적으로 낮추어 878.384점이라는 압도적 고득점을 기록했습니다! 2위는 '베이스라인 1 지도학습 신경망 역산(MLP)'입니다: 3차 포화를 정확히 예측하여 4.0스텝 만에 안착하며 779.169점을 얻었습니다. 3위는 '베이스라인 0 고전 선형 SVD(α=0.45)'입니다: 안정적이지만 9.4스텝이 소요되어 속도 점수가 깎이며 635.277점에 머물렀습니다. 4위는 '트랙 2 기준 RL 에이전트'입니다: 잔차 가중치가 아직 미학습 상태이므로 단순 감쇠 SVD로 동작하여 10.8스텝, 590.559점을 기록했습니다. 우측 궤적 그래프를 보시면 BO와 MLP는 4스텝 이내에 0.030 mm 목표선 아래로 수직 낙하하는 반면, SVD와 RL은 완만한 계단식 하강을 보입니다.
\end{tcolorbox}
\vspace{0.2em}
\begin{tcolorbox}[colback=lightbg,colframe=bordergray,left=6pt,right=6pt,top=3pt,bottom=3pt,boxrule=0.6pt,arc=2.5pt]
\small
\noindent\textbf{\color{navy}교수학적 강조점 (Focus Point):}\enspace 590점에서 878점에 이르는 점수 격차를 짚어주며, 참가자들의 튜닝에 따라 순위가 실시간으로 뒤바뀔 수 있음을 강조합니다.\par\vspace{0.3em}
\noindent\textbf{\color{successgreen}예상 질문 및 모범 답변 (Q\&A):}\enspace Q: RL(트랙 2)이 BO(트랙 3)를 이길 수 있나요? A: 당연히 가능합니다! 주최 측 솔루션에서 잔차를 제대로 학습시킨 RL 에이전트는 무려 939.981점으로 전체 1위를 차지했습니다.\par\vspace{0.3em}
\noindent\textbf{\color{accent}다음 슬라이드 연결 멘트 (Transition):}\enspace 각 제어 기법들의 장단점, 연산 지연시간, 샘플 효율성을 종합 비교 분석해보겠습니다.
\end{tcolorbox}
\clearpage
\hypertarget{slide.33}{}
\phantomsection
\addcontentsline{toc}{section}{슬라이드 33: 제어 방식별 종합 비교 요약 (기법, 연산 지연 및 성능)}
\noindent\begin{tcolorbox}[colback=navy,colframe=navy,left=6pt,right=6pt,top=3pt,bottom=3pt,boxrule=0pt,arc=2pt]
{\large\bfseries\color{white}슬라이드 33:\enspace 제어 방식별 종합 비교 요약 (기법, 연산 지연 및 성능)}\hfill{\footnotesize\bfseries\color{white!85!lightbg}Method Comparison Summary \enspace$|$\enspace 슬라이드 33 / 35}
\end{tcolorbox}
\vspace{0.2em}
\begin{tcolorbox}[colback=lightbg,colframe=bordergray,title={\footnotesize\bfseries\color{darkgray}슬라이드 시각 구성 요약},fontupper=\small,left=6pt,right=6pt,top=2.5pt,bottom=2.5pt,boxrule=0.6pt,arc=2.5pt]
종합 비교 매트릭스 표: 기법별 핵심 메커니즘, 필요 데이터량, 1스텝당 추론 지연시간, 장단점, 기준 점수 및 주최 측 솔루션 점수 비교.
\end{tcolorbox}
\vspace{0.2em}
\begin{tcolorbox}[colback=white,colframe=primaryblue,title={\footnotesize\bfseries\color{white}발표자 구두 대본 (Spoken Script)},fontupper=\small,left=6pt,right=6pt,top=3pt,bottom=3pt,boxrule=0.8pt,arc=2.5pt]
이 표는 오늘 다룬 모든 제어 기법들의 장단점과 핵심 특성을 집대성한 비교 매트릭스입니다. 고전 SVD는 사전 훈련 데이터가 전혀 필요 없고 약 0.01 밀리초로 매우 빠르지만, 비선형 포화를 전혀 다루지 못합니다. 지도학습 MLP 역산은 1,500개 오프라인 샷으로 비선형성을 완벽히 포착하여 기본 779점, 주최 측 솔루션에서는 924점을 달성합니다. 다만 매 스텝 최적화 수치해를 구해야 하므로 스텝당 약 27 밀리초가 소요됩니다. 트랙 2 강화학습 정책은 물리 사전 지식과 잔차를 결합한 직접 피드백 정책이므로 스텝당 약 0.01 밀리초로, MLP 역산보다 천 배 이상 빠릅니다. 기준 점수는 590점이지만, 선형 사전 지식을 정확한 포화 모델로 교체하면 940점을 기록합니다. 트랙 3 베이지안 최적화는 단 300개의 캘리브레이션 샷만 필요한 가장 샘플 효율적인 방법으로, 기준 설정에서 878점, 주최 측 솔루션에서 910점을 기록합니다. 다만 스텝당 약 66 밀리초로 네 가지 중 가장 느립니다. 요약하자면, 초고속 실시간성이 최우선이라면 RL, 빔 측정 샷 수가 극도로 제한된다면 BO, 풍부한 시뮬레이션 데이터가 있다면 지도학습 대리 모델이 정답입니다.
\end{tcolorbox}
\vspace{0.2em}
\begin{tcolorbox}[colback=lightbg,colframe=bordergray,left=6pt,right=6pt,top=3pt,bottom=3pt,boxrule=0.6pt,arc=2.5pt]
\small
\noindent\textbf{\color{navy}교수학적 강조점 (Focus Point):}\enspace 연산 속도, 필요 데이터량, 알고리즘 복잡도 간의 삼각 트레이드오프를 짚어주며 강의 전체를 유기적으로 정리합니다.\par\vspace{0.3em}
\noindent\textbf{\color{successgreen}예상 질문 및 모범 답변 (Q\&A):}\enspace Q: 초심자에게는 어떤 트랙을 가장 추천하시나요? A: 딥러닝 입문자에게는 직관적인 트랙 1(TODO 1)을 먼저 권장하며, 이후 트랙 3이나 트랙 2로 확장하는 것이 좋습니다.\par\vspace{0.3em}
\noindent\textbf{\color{accent}다음 슬라이드 연결 멘트 (Transition):}\enspace 그렇다면 오늘 오후 해커톤에서 리더보드 상위권에 오르기 위한 구체적인 실전 공략 전략을 공유해 드리겠습니다.
\end{tcolorbox}
\clearpage
\hypertarget{slide.34}{}
\phantomsection
\addcontentsline{toc}{section}{슬라이드 34: 해커톤 참가자를 위한 실전 공략 전략 (TODO 1, 2, 3)}
\noindent\begin{tcolorbox}[colback=navy,colframe=navy,left=6pt,right=6pt,top=3pt,bottom=3pt,boxrule=0pt,arc=2pt]
{\large\bfseries\color{white}슬라이드 34:\enspace 해커톤 참가자를 위한 실전 공략 전략 (TODO 1, 2, 3)}\hfill{\footnotesize\bfseries\color{white!85!lightbg}Strategic Advice \enspace$|$\enspace 슬라이드 34 / 35}
\end{tcolorbox}
\vspace{0.2em}
\begin{tcolorbox}[colback=lightbg,colframe=bordergray,title={\footnotesize\bfseries\color{darkgray}슬라이드 시각 구성 요약},fontupper=\small,left=6pt,right=6pt,top=2.5pt,bottom=2.5pt,boxrule=0.6pt,arc=2.5pt]
트랙별 핵심 튜닝 전략 카드: 트랙 1(TODO 1 대리 모델 고도화), 트랙 2(TODO 2 잔차 정책 학습 및 감쇠 튜닝), 트랙 3(TODO 3 온라인 GP 갱신 및 동적 κ 스케줄).
\end{tcolorbox}
\vspace{0.2em}
\begin{tcolorbox}[colback=white,colframe=primaryblue,title={\footnotesize\bfseries\color{white}발표자 구두 대본 (Spoken Script)},fontupper=\small,left=6pt,right=6pt,top=3pt,bottom=3pt,boxrule=0.8pt,arc=2.5pt]
오늘 오후 미니 해커톤에서 승리하기 위한 실전 튜닝 전략 가이드입니다. 트랙 1(TODO 1 대리 모델 튜닝) 참가자 여러분: 기본 제공된 MLP 설정에 머무르지 마십시오! tanh 같은 매끄러운 활성화 함수를 시도해보고, 오프라인 학습 시 solver='lbfgs'를 적용해보십시오. 은닉층 구조를 (64, 32)로 경량화하거나 (256, 128)로 확장하고 L2 규제를 조절하십시오. 테스트 RMSE를 기본 설정의 0.109 mm에서 약 0.012 mm로 낮추면 베이스라인 1의 점수가 779점에서 약 900점으로 올라갑니다. 트랙 2(TODO 2 RL 정책 튜닝) 참가자 여러분: 첫째, 물리 사전 지식의 선형 게인 α를 0.8~1.2 사이로 최적화하십시오. 둘째, 감쇠 계수 κ를 0~0.05 사이에서 시험해 보십시오. 게인에 따라 도움이 되기도 하고 손해가 되기도 합니다. 셋째, info['currents']의 누적 전류를 이용해 자기 포화를 보상하십시오. 이것만으로 924점에 도달합니다. 또는 잔차 행렬 W\_learned가 같은 역할을 하도록 학습시키십시오. 트랙 3(TODO 3 베이지안 최적화) 참가자 여러분: κ 값을 0.0~0.5 사이에서 테스트하고, step\_callback을 활용해 측정 데이터를 온라인으로 GP에 반영하십시오. 초기에는 탐색하고 후기에는 활용하는 동적 κ 스케줄러도 훌륭한 전략입니다. 무엇보다 가장 중요한 황금률은 '강건성이 최우선'이라는 점입니다! 단 한 번의 시드 실패(0점)도 없도록 안정성을 먼저 확보한 후 속도와 정밀도를 올리십시오.
\end{tcolorbox}
\vspace{0.2em}
\begin{tcolorbox}[colback=lightbg,colframe=bordergray,left=6pt,right=6pt,top=3pt,bottom=3pt,boxrule=0.6pt,arc=2.5pt]
\small
\noindent\textbf{\color{navy}교수학적 강조점 (Focus Point):}\enspace 환경의 물리 상수나 평가 시드를 조작하지 않고 정정당당하게 알고리즘 자체를 개선해야 한다는 규칙을 환기하며 도전 의식을 북돋웁니다.\par\vspace{0.3em}
\noindent\textbf{\color{successgreen}예상 질문 및 모범 답변 (Q\&A):}\enspace Q: 두 가지 방식을 융합해도 되나요? A: 물론입니다! GP 예측을 사전 지식으로 활용하는 하이브리드 제어기는 실제 연구에서도 매우 강력한 접근법입니다.\par\vspace{0.3em}
\noindent\textbf{\color{accent}다음 슬라이드 연결 멘트 (Transition):}\enspace 이제 오늘 강의를 마무리하며 핵심 교훈을 정리하고, Thorsten Hellert 박사님의 제2강으로 넘어가겠습니다.
\end{tcolorbox}
\clearpage
\hypertarget{slide.35}{}
\phantomsection
\addcontentsline{toc}{section}{슬라이드 35: 강의 요약, 핵심 교훈 및 제2강 연계 안내}
\noindent\begin{tcolorbox}[colback=navy,colframe=navy,left=6pt,right=6pt,top=3pt,bottom=3pt,boxrule=0pt,arc=2pt]
{\large\bfseries\color{white}슬라이드 35:\enspace 강의 요약, 핵심 교훈 및 제2강 연계 안내}\hfill{\footnotesize\bfseries\color{white!85!lightbg}Summary \& Transition \enspace$|$\enspace 슬라이드 35 / 35}
\end{tcolorbox}
\vspace{0.2em}
\begin{tcolorbox}[colback=lightbg,colframe=bordergray,title={\footnotesize\bfseries\color{darkgray}슬라이드 시각 구성 요약},fontupper=\small,left=6pt,right=6pt,top=2.5pt,bottom=2.5pt,boxrule=0.6pt,arc=2.5pt]
핵심 교훈 정리 카드: 물리 지식과 ML의 융합 필수성, 비선형 포화 극복, 이어서 진행될 제2강 '가속기 과학을 위한 에이전틱 AI(Thorsten Hellert 박사)' 안내 배너.
\end{tcolorbox}
\vspace{0.2em}
\begin{tcolorbox}[colback=white,colframe=primaryblue,title={\footnotesize\bfseries\color{white}발표자 구두 대본 (Spoken Script)},fontupper=\small,left=6pt,right=6pt,top=3pt,bottom=3pt,boxrule=0.8pt,arc=2.5pt]
제1강을 마치며 핵심 교훈을 정리하겠습니다. 첫째, '가속기 물리학과 머신러닝의 융합'은 선택이 아닌 필수입니다. 순수 블랙박스 신경망은 외삽 영역에서 물리 법칙을 위반하며, 순수 선형 광학 모델은 자기 포화라는 비선형 현실 앞에서 무력해집니다. 물리 사전 지식에 딥러닝 대리 모델, 적응형 RL, 베이지안 최적화를 결합할 때 비로소 센서 노이즈 바닥까지 도달하는 강건한 궤도 제어가 완성됩니다. 둘째, 해커톤 채점 공식은 10 ㎛ 노이즈 바닥까지의 강건성, 속도, 정밀도를 종합 평가하도록 설계되었습니다. 이제 잠시 휴식 후, 미국 LBNL의 Thorsten Hellert 박사님께서 진행하시는 제2강 '가속기 연구를 위한 에이전틱 AI'가 이어집니다. 최신 거대언어모델 에이전트가 가속기 제어 환경을 자율적으로 분석하고 시뮬레이션을 오케스트레이션하는 놀라운 최신 기술을 만나보실 수 있습니다. 경청해 주셔서 대단히 감사드리며, 오후 해커톤에서 여러분의 멋진 활약을 기대하겠습니다!
\end{tcolorbox}
\vspace{0.2em}
\begin{tcolorbox}[colback=lightbg,colframe=bordergray,left=6pt,right=6pt,top=3pt,bottom=3pt,boxrule=0.6pt,arc=2.5pt]
\small
\noindent\textbf{\color{navy}교수학적 강조점 (Focus Point):}\enspace 열정적인 어조로 강의를 마무리하고, 질의응답 시간을 가진 뒤 제2강 Thorsten Hellert 박사에게 마이크를 넘깁니다.\par\vspace{0.3em}
\noindent\textbf{\color{successgreen}예상 질문 및 모범 답변 (Q\&A):}\enspace Q: 강의 슬라이드와 강의 노트는 어디서 볼 수 있나요? A: 저장소 루트의 presentations/ 및 lecture\_notes/ 디렉토리에 PDF와 마크다운으로 모두 제공됩니다.\par\vspace{0.3em}
\noindent\textbf{\color{accent}다음 슬라이드 연결 멘트 (Transition):}\enspace 감사합니다. 이어서 제2강과 오후 해커톤을 유익하게 즐겨주시기 바랍니다!
\end{tcolorbox}
\end{document}